\chapter{Singular Homology and Cohomology}\label{singular-homology}
\chapterstatus{complete}

\section{Introduction}\label{singular-homology:section-introduction}

Singular homology assigns to every space $X$ abelian groups $H_n(X)$, functorial in $X$ and invariant under homotopy. Unlike
homotopy groups, they are computable: the long exact sequence of a pair, excision
(Theorem~\ref{singular-homology:theorem-excision}) and the Mayer--Vietoris sequence reduce computations to simple pieces, and for CW
complexes cellular homology (Theorem~\ref{singular-homology:theorem-cellular}) computes them from the cells. The last section compares homology with homotopy
groups: $H_1$ is the abelianised fundamental group, and the first nonvanishing homotopy group of a simply connected space is its first nonvanishing homology group
(Theorem~\ref{singular-homology:theorem-hurewicz-general}). We follow
\cite[Chapter 2]{hatcher-at}; see also \cite{spanier-at}.

\noindent\textbf{Prerequisites.} Chapter~\ref{homotopy} (Homotopy and Covering Spaces) and Chapter~\ref{homological-algebra}
(Homological Algebra).

Abelian groups are written additively. Chain complexes are as in Notation~\ref{homological-algebra:notation-chain-complexes}, with
differential $\partial$ of degree $-1$.

\section{Singular chains}\label{singular-homology:section-chains}

\begin{definition}\label{singular-homology:definition-simplex}
The \emph{standard $n$-simplex} is $\Delta^n = \{(t_0, \ldots, t_n) \in \RR^{n+1} : t_i \geq 0, \sum_i t_i = 1\}$, with vertices the standard basis
vectors $e_0, \ldots, e_n$. For points $w_0, \ldots, w_k$ of some $\RR^N$ we write $[w_0, \ldots, w_k] \colon \Delta^k \to \RR^N$ for the affine map
$(t_i) \mapsto \sum_i t_iw_i$. The \emph{face maps} are $\delta^i = [e_0, \ldots, \widehat{e_i}, \ldots, e_n] \colon \Delta^{n-1} \to \Delta^n$, where the hat
means omission.
\end{definition}

\begin{definition}\label{singular-homology:definition-singular-homology}
A \emph{singular $n$-simplex} in a space $X$ is a map $\sigma \colon \Delta^n \to X$. The group $C_n(X)$ of \emph{singular $n$-chains} is the free abelian group
(Definition~\ref{groups:definition-free-abelian}) with basis the singular $n$-simplices; its elements are finite formal sums $\sum_i n_i\sigma_i$ with
$n_i \in \ZZ$. For $n < 0$, $C_n(X) = 0$. The \emph{boundary} $\partial \colon C_n(X) \to C_{n-1}(X)$ is the homomorphism with
\[
\partial\sigma = \sum_{i=0}^n (-1)^i \sigma \circ \delta^i.
\]
The \emph{singular homology} of $X$ is $H_n(X) = \ker(\partial \colon C_n \to C_{n-1})/\im(\partial \colon C_{n+1} \to C_n)$. A map $f \colon X \to Y$ induces
$f_\# \colon C_n(X) \to C_n(Y)$, $\sigma \mapsto f \circ \sigma$, and $f_* \colon H_n(X) \to H_n(Y)$.
\end{definition}

\emph{What the definition says in low dimensions.} A singular $0$-simplex is a point of $X$. A singular $1$-simplex is a path, via $\Delta^1 \cong I$, $(t_0, t_1) \mapsto t_1$,
and $\partial\sigma = \sigma(e_1) - \sigma(e_0)$ is ``endpoint minus starting point''. A singular $2$-simplex $\sigma$ is a map of a triangle, with three edges
$\tau_i = \sigma \circ \delta^i$: $\tau_2$ runs from $\sigma(e_0)$ to $\sigma(e_1)$, $\tau_0$ from $\sigma(e_1)$ to $\sigma(e_2)$, and $\tau_1$ from $\sigma(e_0)$ to $\sigma(e_2)$.
So $\partial\sigma = \tau_2 + \tau_0 - \tau_1$ is the boundary of the triangle traversed once: along $\tau_2$, then $\tau_0$, then back along $\tau_1$. The signs in the
definition of $\partial$ are chosen so that this works in every dimension, and so that $\partial\partial = 0$: the boundary of a boundary is empty.

\begin{example}\label{singular-homology:example-cycles}
\begin{enumerate}
\item A loop $\gamma$, regarded as a $1$-simplex, is a cycle: $\partial\gamma = \gamma(1) - \gamma(0) = 0$. More generally, if $f_1, \ldots, f_k$ are paths such that each
$f_i$ ends where $f_{i+1}$ starts and $f_k$ ends where $f_1$ starts, then $f_1 + \cdots + f_k$ is a cycle, a ``closed polygon'' of paths, since the endpoints cancel
in pairs.
\item The constant $2$-simplex $c_2$ at a point $x$ has all three edges equal to the constant $1$-simplex $c_1$, so $\partial c_2 = c_1 - c_1 + c_1 = c_1$: a constant
path is a boundary, and is zero in $H_1(X)$.
\item The $1$-cycles that are boundaries are those that ``bound a surface'' made of triangles. So $H_1(X)$ counts closed curves in $X$ that are not filled in by
surfaces, with multiplicities and signs. In the circle, the cycle formed by the upper and the lower half circle is not a boundary
(Example~\ref{singular-homology:example-circle-generator} below), while in a disc every $1$-cycle is a boundary
(Theorem~\ref{singular-homology:theorem-homotopy-invariance} below). The precise relation with closed curves up to homotopy is
Theorem~\ref{singular-homology:theorem-hurewicz}.
\end{enumerate}
\end{example}

\begin{lemma}\label{singular-homology:lemma-boundary-squared}
$\partial\partial = 0$, and $f_\#$ is a morphism of chain complexes. Hence $H_n$ is a functor $\cat{Top} \to \cat{Ab}$.
\end{lemma}

\begin{proof}
For $i < j$ one has $\delta^j\delta^i = \delta^i\delta^{j-1}$ as maps $\Delta^{n-2} \to \Delta^n$: both are the affine map omitting $e_i$ and $e_j$. So
\[
\partial\partial\sigma = \sum_{j}\sum_{i} (-1)^{i+j}\sigma\delta^j\delta^i
= \sum_{i < j} (-1)^{i+j}\sigma\delta^i\delta^{j-1} + \sum_{i \geq j}(-1)^{i+j}\sigma\delta^j\delta^i,
\]
and substituting $(i, j - 1) \mapsto (j', i')$ in the first sum shows that it is the negative of the second. Clearly $\partial f_\# = f_\#\partial$.
\end{proof}

\begin{proposition}\label{singular-homology:proposition-h0-point}
\begin{enumerate}
\item If $(X_\alpha)$ are the path components of $X$, then $H_n(X) \cong \bigoplus_\alpha H_n(X_\alpha)$.
\item If $X$ is path-connected and nonempty, then $H_0(X) \cong \ZZ$, via $\sum n_ix_i \mapsto \sum n_i$.
\item For a point, $H_0 = \ZZ$ and $H_n = 0$ for $n \neq 0$.
\end{enumerate}
\end{proposition}

\begin{proof}
(1) The image of a singular simplex is path-connected, so it lies in one $X_\alpha$, and $C_\bullet(X) = \bigoplus_\alpha C_\bullet(X_\alpha)$ as chain complexes.
(2) Singular $0$-simplices are points and $\partial\sigma = \sigma(e_1) - \sigma(e_0)$ for a $1$-simplex, a path. So the \emph{augmentation}
$\varepsilon(\sum n_ix_i) = \sum n_i$ vanishes on boundaries and is surjective. If $\varepsilon(\sum n_ix_i) = 0$, fix $x_0$ and paths $\gamma_i$ from $x_0$
to $x_i$; then $\sum n_ix_i = \sum n_i(x_i - x_0) = \partial(\sum n_i\gamma_i)$. (3) There is exactly one singular $n$-simplex $\sigma_n$ for each $n$, and
$\partial\sigma_n = \sum_{i=0}^n(-1)^i\sigma_{n-1}$ is $0$ for $n$ odd and $\sigma_{n-1}$ for $n \geq 2$ even. So the complex is
$\cdots \to \ZZ \xrightarrow{1} \ZZ \xrightarrow{0} \ZZ \xrightarrow{1} \ZZ \xrightarrow{0} \ZZ \to 0$, with homology $\ZZ$ in degree $0$ only.
\end{proof}

\begin{definition}\label{singular-homology:definition-reduced}
The \emph{reduced homology} $\tilde{H}_n(X)$ of a nonempty space is the homology of the augmented complex
$\cdots \to C_1(X) \to C_0(X) \xrightarrow{\varepsilon} \ZZ \to 0$, with $\ZZ$ in degree $-1$. So $\tilde{H}_n(X) = H_n(X)$ for $n > 0$,
$H_0(X) \cong \tilde{H}_0(X) \oplus \ZZ$, and $\tilde{H}_n(\mathrm{pt}) = 0$ for all $n$.
\end{definition}

\section{Homotopy invariance}\label{singular-homology:section-homotopy-invariance}

\begin{theorem}\label{singular-homology:theorem-homotopy-invariance}
If $f \simeq g \colon X \to Y$, then $f_\#$ and $g_\#$ are chain homotopic, so $f_* = g_* \colon H_n(X) \to H_n(Y)$. Consequently homotopy equivalences
induce isomorphisms, and contractible spaces have $\tilde{H}_\bullet = 0$.
\end{theorem}

\begin{proof}
Let $F \colon X \times I \to Y$ be a homotopy from $f$ to $g$. In $\Delta^n \times I \subseteq \RR^{n+2}$ let $v_i = (e_i, 0)$ and $w_i = (e_i, 1)$. Define the
\emph{prism operator} $P \colon C_n(X) \to C_{n+1}(Y)$ by
\[
P(\sigma) = \sum_{i=0}^n (-1)^i F \circ (\sigma \times \id_I) \circ [v_0, \ldots, v_i, w_i, \ldots, w_n].
\]
Write $[\ldots]$ for $F \circ (\sigma \times \id) \circ [\ldots]$. Then
\[
\partial P(\sigma) = \sum_{j \leq i}(-1)^{i+j}[v_0, \ldots, \widehat{v_j}, \ldots, v_i, w_i, \ldots, w_n]
+ \sum_{j \geq i}(-1)^{i+j+1}[v_0, \ldots, v_i, w_i, \ldots, \widehat{w_j}, \ldots, w_n].
\]
In the terms with $i = j$, the term $[v_0, \ldots, v_{i-1}, w_i, \ldots, w_n]$ of the first sum cancels the term with index $i - 1$ of the second sum, for
$1 \leq i \leq n$. What remains of them is $[w_0, \ldots, w_n] = g \circ \sigma$ (first sum, $i = 0$) and $-[v_0, \ldots, v_n] = -f \circ \sigma$ (second sum,
$i = n$). Since $\sigma \circ \delta^j$ corresponds to the vertices with index $j$ omitted,
\[
P(\partial\sigma) = \sum_{i < j}(-1)^{i+j}[v_0, \ldots, v_i, w_i, \ldots, \widehat{w_j}, \ldots, w_n] + \sum_{i > j}(-1)^{i-1+j}[v_0, \ldots, \widehat{v_j}, \ldots, v_i, w_i, \ldots, w_n],
\]
which is the negative of the terms with $i \neq j$ above. Hence $\partial P + P\partial = g_\# - f_\#$, and chain homotopic maps induce the same map on homology
(Lemma~\ref{homological-algebra:lemma-homotopy-invariance}). The rest follows from functoriality and Proposition~\ref{singular-homology:proposition-h0-point}(3).
\end{proof}

\begin{remark}\label{singular-homology:remark-prism}
\emph{What the prism operator does.} A homotopy $F$ sweeps each simplex $\sigma$ of $X$ through the \emph{prism} $\Delta^n \times I$, and $P(\sigma)$ is this prism, pushed into $Y$ by $F \circ (\sigma \times \id)$.
The prism is not a simplex, and singular chains are made of simplices, so it has to be cut into simplices first: the $n + 1$ affine simplices $[v_0, \ldots, v_i, w_i, \ldots, w_n]$, $0 \leq i \leq n$, fill
$\Delta^n \times I$ and meet along common faces, and the signs $(-1)^i$ orient them coherently. The identity $\partial P + P\partial = g_\# - f_\#$ is then the geometric statement that the boundary of the prism
consists of its top $\Delta^n \times \{1\}$, its bottom $\Delta^n \times \{0\}$ with the opposite orientation, and the prisms over the faces of $\Delta^n$, the sides.
\begin{enumerate}
\item For $n = 0$, $\sigma$ is a point $x$, the prism is the segment $\{x\} \times I$, and $P(x)$ is the path $t \mapsto F(x, t)$ from $f(x)$ to $g(x)$. So $\partial P(x) = g(x) - f(x)$, while $P(\partial x) = 0$.
\item For $n = 1$, $\sigma$ is a path from $x_0$ to $x_1$, and $\Delta^1 \times I$ is a square with bottom edge $[v_0, v_1]$, top edge $[w_0, w_1]$ and vertical sides $[v_0, w_0]$ and $[v_1, w_1]$. The diagonal from $v_0$ to $w_1$
cuts it into the two triangles of $P(\sigma) = [v_0, w_0, w_1] - [v_0, v_1, w_1]$. Their boundaries are
\[
\partial[v_0, w_0, w_1] = [w_0, w_1] - [v_0, w_1] + [v_0, w_0], \qquad \partial[v_0, v_1, w_1] = [v_1, w_1] - [v_0, w_1] + [v_0, v_1],
\]
and in the difference the diagonal $[v_0, w_1]$ cancels:
\[
\partial P(\sigma) = [w_0, w_1] - [v_0, v_1] + [v_0, w_0] - [v_1, w_1] = g_\#\sigma - f_\#\sigma - P(\partial\sigma),
\]
since $P(\partial\sigma) = P(x_1) - P(x_0)$ is the path of $x_1$ minus the path of $x_0$, that is, $[v_1, w_1] - [v_0, w_0]$. This is the picture behind the general computation: the interior cuts cancel in pairs, and what is
left is the top, minus the bottom, minus the sides.
\end{enumerate}
In the de Rham complex the same role is played by the homotopy operator $K$ of Chapter~\ref{de-rham}, which integrates a form over the $I$-direction. There no subdivision is needed, because forms can be
integrated over the product $\Delta^n \times I$ directly; for chains, which are sums of simplices, the product has to be triangulated, and this is the only reason why $P$ looks more complicated than $K$.
\end{remark}

\section{Relative homology and exact sequences}\label{singular-homology:section-relative}

\begin{definition}\label{singular-homology:definition-relative}
For a subspace $A \subseteq X$, $C_n(X, A) = C_n(X)/C_n(A)$ with the induced boundary, and $H_n(X, A)$ is its homology, the \emph{relative homology}. A
map of pairs $f \colon (X, A) \to (Y, B)$ (a map with $f(A) \subseteq B$) induces $f_* \colon H_n(X, A) \to H_n(Y, B)$, and homotopic maps of pairs (through
maps of pairs) induce the same maps, since the prism operator preserves chains in $A$.
\end{definition}

\begin{theorem}\label{singular-homology:theorem-long-exact-sequence}
For $B \subseteq A \subseteq X$ there is a natural long exact sequence
\[
\cdots \to H_n(A, B) \to H_n(X, B) \to H_n(X, A) \xrightarrow{\partial} H_{n-1}(A, B) \to \cdots \to H_0(X, A) \to 0.
\]
For $B = \emptyset$ this is the long exact sequence of the pair $(X, A)$, where $\partial[c] = [\partial c]$ for a relative cycle $c$. The same holds with
reduced homology when $B = \emptyset$ and $A \neq \emptyset$.
\end{theorem}

\begin{proof}
The sequence $0 \to C(A, B) \to C(X, B) \to C(X, A) \to 0$ of chain complexes is exact. Apply Theorem~\ref{homological-algebra:theorem-long-exact-sequence}. For
reduced homology use the augmented complexes, whose quotient is still $C(X, A)$.
\end{proof}

\begin{example}\label{singular-homology:example-pair-point}
For $x_0 \in X$, the long exact sequence of $(X, x_0)$ in reduced homology gives $H_n(X, x_0) \cong \tilde{H}_n(X)$. If $A$ is a deformation retract of $X$,
then $H_n(X, A) = 0$ for all $n$, by exactness of the long exact sequence, since $H_n(A) \to H_n(X)$ is an isomorphism for all $n$
(Theorem~\ref{singular-homology:theorem-homotopy-invariance}).
\end{example}

\begin{example}\label{singular-homology:example-interval}
\emph{What relative cycles are.} A chain $c \in C_n(X)$ represents a class in $H_n(X, A)$ when $\partial c \in C_{n-1}(A)$: it is closed ``up to $A$''. The simplest case is
$X = I = [0, 1]$, $A = \partial I = \{0, 1\}$. Since $I$ is contractible, $\tilde{H}_\bullet(I) = 0$, and the reduced long exact sequence
$\tilde{H}_n(I) \to H_n(I, \partial I) \xrightarrow{\partial} \tilde{H}_{n-1}(\partial I) \to \tilde{H}_{n-1}(I)$ shows that $\partial$ is an isomorphism for every $n$. By
Proposition~\ref{singular-homology:proposition-h0-point}, $\tilde{H}_\bullet(\partial I)$ is zero except in degree $0$, where it is the kernel of the augmentation
$\ZZ\{0\} \oplus \ZZ\{1\} \to \ZZ$, generated by $\{1\} - \{0\}$. So $H_n(I, \partial I)$ is $\ZZ$ for $n = 1$ and $0$ otherwise. A generator is the identity path
$\iota(t) = t$: it is not a cycle in $I$, but its boundary $\{1\} - \{0\}$ lies in $\partial I$, and $\partial[\iota] = [\{1\} - \{0\}]$ is the generator. This is the
one-dimensional case of $H_n(D^n, S^{n-1}) \cong \ZZ$ (Proposition~\ref{singular-homology:proposition-good-pair} below).
\end{example}

\section{Excision and Mayer--Vietoris}\label{singular-homology:section-excision}

\begin{definition}\label{singular-homology:definition-small-chains}
Let $\mathcal{U}$ be a family of subspaces of $X$ whose interiors cover $X$. Then $C^{\mathcal{U}}_n(X) \subseteq C_n(X)$ is the subgroup generated by singular
simplices with image in some $U \in \mathcal{U}$. It is a subcomplex.
\end{definition}

\begin{lemma}\label{singular-homology:lemma-subdivision}
There are homomorphisms $S \colon C_n(X) \to C_n(X)$ and $T \colon C_n(X) \to C_{n+1}(X)$, natural in $X$ ($f_\#S = Sf_\#$, $f_\#T = Tf_\#$), with
$\partial S = S\partial$ and $\partial T + T\partial = 1 - S$, such that $S(\sigma)$ and $T(\sigma)$ are combinations of simplices with image in $\sigma(\Delta^n)$.
Moreover $S^m(\sigma) = \sigma_\#(S^m\iota_n)$ with $\iota_n = [e_0, \ldots, e_n]$, and $S^m\iota_n$ is a combination of affine simplices of diameter at most
$(n/(n+1))^m\sqrt{2}$.
\end{lemma}

\begin{proof}
\emph{Linear chains.} For a convex $Y \subseteq \RR^N$ let $LC_n(Y) \subseteq C_n(Y)$ be generated by the affine simplices $[w_0, \ldots, w_n]$; it is a subcomplex. For
$b \in Y$ the \emph{cone} $b \colon LC_n(Y) \to LC_{n+1}(Y)$, $[w_0, \ldots, w_n] \mapsto [b, w_0, \ldots, w_n]$, satisfies $\partial b(\lambda) = \lambda - b(\partial\lambda)$
for $n \geq 1$, and $\partial b(\lambda) = \lambda - \varepsilon(\lambda)[b]$ for $n = 0$, by the formula for $\partial$. For an affine simplex $\lambda$ let $b_\lambda$ be
its barycentre $\frac{1}{n+1}\sum_i w_i$. Define $S = \id$ on $LC_0$ and $S(\lambda) = b_\lambda(S\partial\lambda)$ for $n \geq 1$; define $T = 0$ on $LC_0$ and
$T(\lambda) = b_\lambda(\lambda - T\partial\lambda)$ for $n \geq 1$.

By induction on $n$: $\partial S\lambda = S\partial\lambda - b_\lambda(\partial S\partial\lambda) = S\partial\lambda - b_\lambda(S\partial\partial\lambda) = S\partial\lambda$ for $n \geq 2$,
and for $n = 1$, $\partial S\lambda = S\partial\lambda - \varepsilon(\partial\lambda)[b_\lambda] = S\partial\lambda$. Similarly, using $\partial T\partial\lambda = \partial\lambda - S\partial\lambda$ by
induction, $\partial T\lambda = \lambda - T\partial\lambda - b_\lambda(\partial\lambda - \partial T\partial\lambda) = \lambda - T\partial\lambda - b_\lambda(S\partial\lambda) = \lambda - T\partial\lambda - S\lambda$.
Affine maps preserve barycentres and cones, so $S$ and $T$ commute with $g_\#$ for affine maps $g$.

\emph{General chains.} Put $S\sigma = \sigma_\#S\iota_n$ and $T\sigma = \sigma_\#T\iota_n$. Naturality is clear, and since $\sigma \circ \delta^i = \sigma_\#\delta^i$ with $\delta^i$ affine,
$\partial S\sigma = \sigma_\#S\partial\iota_n = \sum_i(-1)^i\sigma_\#\delta^i_\#S\iota_{n-1} = S\partial\sigma$, and likewise for $T$. By naturality
$S^m\sigma = \sigma_\#S^m\iota_n$.

\emph{Diameters.} The diameter of an affine simplex is the largest distance between its vertices, since $|\sum_i t_iw_i - y| \leq \max_i|w_i - y|$. The simplices of
$S\lambda$ are of the form $[b_\lambda, u_0, \ldots, u_{n-1}]$ with $[u_0, \ldots, u_{n-1}]$ a simplex of $S$ applied to a face of $\lambda$. For a vertex $w_i$,
$|b_\lambda - w_i| = |\sum_j (w_j - w_i)|/(n+1) \leq \frac{n}{n+1}\operatorname{diam}\lambda$, so all points of $\lambda$ are within this distance of $b_\lambda$;
and by induction the simplices of $S$ of a face have diameter at most $\frac{n-1}{n}\operatorname{diam}\lambda$. So the simplices of $S\lambda$ have diameter at most
$\frac{n}{n+1}\operatorname{diam}\lambda$. Iterate, using $\operatorname{diam}\Delta^n = \sqrt{2}$.
\end{proof}

\begin{proposition}\label{singular-homology:proposition-small-chains}
The inclusion $\iota \colon C^{\mathcal{U}}(X) \to C(X)$ is a chain homotopy equivalence. Hence it induces isomorphisms on homology.
\end{proposition}

\begin{proof}
For a singular simplex $\sigma \colon \Delta^n \to X$ let $m(\sigma)$ be the least $m \geq 0$ such that every subset of $\Delta^n$ of diameter at most $(n/(n+1))^m\sqrt2$ lies
in $\sigma^{-1}(\operatorname{int}U)$ for some $U \in \mathcal{U}$. It exists by the Lebesgue number lemma (Lemma~\ref{general-topology:lemma-lebesgue-number}), and
by Lemma~\ref{singular-homology:lemma-subdivision}, $S^i\sigma \in C^{\mathcal{U}}(X)$ for $i \geq m(\sigma)$. The face maps are isometric embeddings and
$\frac{n-1}{n} < \frac{n}{n+1}$, so $m(\sigma \circ \delta^j) \leq m(\sigma)$.

Let $D_m = \sum_{0 \leq i < m}TS^i$, so that $\partial D_m + D_m\partial = \sum_i(1 - S)S^i = 1 - S^m$. Define $D\sigma = D_{m(\sigma)}\sigma$ and $\rho = 1 - \partial D - D\partial$. Then
$\rho$ is a chain map, and
\[
\rho(\sigma) = S^{m(\sigma)}\sigma + (D_{m(\sigma)} - D)\partial\sigma = S^{m(\sigma)}\sigma + \sum_j (-1)^j\sum_{m(\sigma\delta^j) \leq i < m(\sigma)} TS^i(\sigma\delta^j)
\]
lies in $C^{\mathcal{U}}(X)$, since $T$ preserves $C^{\mathcal{U}}(X)$. So $\rho \colon C(X) \to C^{\mathcal{U}}(X)$ satisfies $1 - \iota\rho = \partial D + D\partial$. As $D$ maps
$C^{\mathcal{U}}(X)$ into itself, the same identity shows $1 - \rho\iota = \partial D + D\partial$ on $C^{\mathcal{U}}(X)$.
\end{proof}

\begin{theorem}[Excision]\label{singular-homology:theorem-excision}
Let $Z \subseteq A \subseteq X$ with $\overline{Z} \subseteq \operatorname{int}A$. Then the inclusion $(X \setminus Z, A \setminus Z) \to (X, A)$ induces isomorphisms on all
$H_n$. Equivalently, if $A, B \subseteq X$ with $\operatorname{int}A \cup \operatorname{int}B = X$, then $H_n(B, A \cap B) \to H_n(X, A)$ is an isomorphism.
\end{theorem}

\begin{proof}
The two statements correspond via $B = X \setminus Z$, since $\operatorname{int}B = X \setminus \overline{Z}$. Let $\mathcal{U} = \{A, B\}$. The map
$C(B)/C(A \cap B) \to C^{\mathcal{U}}(X)/C(A)$ is an isomorphism of complexes: it is surjective as every generator of $C^{\mathcal{U}}$ lies in $C(A)$ or $C(B)$, and its
kernel is $C(B) \cap C(A) = C(A \cap B)$. The map $C^{\mathcal{U}}(X)/C(A) \to C(X)/C(A)$ induces isomorphisms on homology, by
Proposition~\ref{singular-homology:proposition-small-chains} and the five lemma applied to the long exact sequences of $0 \to C(A) \to C^{\mathcal{U}}(X) \to C^{\mathcal{U}}(X)/C(A) \to 0$
and $0 \to C(A) \to C(X) \to C(X, A) \to 0$.
\end{proof}

\begin{theorem}[Mayer--Vietoris]\label{singular-homology:theorem-mayer-vietoris}
If $A, B \subseteq X$ with $\operatorname{int}A \cup \operatorname{int}B = X$, there is a natural long exact sequence
\[
\cdots \to H_n(A \cap B) \xrightarrow{(i_*, -j_*)} H_n(A) \oplus H_n(B) \xrightarrow{k_* + l_*} H_n(X) \xrightarrow{\partial} H_{n-1}(A \cap B) \to \cdots,
\]
where $i, j, k, l$ are the inclusions, and similarly in reduced homology if $A \cap B \neq \emptyset$.
\end{theorem}

\begin{proof}
The sequence $0 \to C(A \cap B) \xrightarrow{(x \mapsto (x, -x))} C(A) \oplus C(B) \xrightarrow{((x, y) \mapsto x + y)} C^{\{A, B\}}(X) \to 0$ is exact: the kernel of the second
map consists of pairs $(x, -x)$ with $x \in C(A) \cap C(B) = C(A \cap B)$. Take the long exact sequence and replace $H_n(C^{\{A,B\}}(X))$ by $H_n(X)$ using
Proposition~\ref{singular-homology:proposition-small-chains}. For reduced homology augment each complex by $\ZZ$, $\ZZ \oplus \ZZ$, $\ZZ$.
\end{proof}

\begin{remark}\label{singular-homology:remark-mayer-vietoris-boundary}
\emph{How to compute the connecting map.} By Proposition~\ref{singular-homology:proposition-small-chains}, every class in $H_n(X)$ is represented by a cycle $z$ that is a sum
of simplices lying in $A$ or in $B$; write $z = a + b$ with $a \in C_n(A)$ and $b \in C_n(B)$ (simplices in $A \cap B$ may go to either side). Then
$\partial a = -\partial b$ lies in $C_{n-1}(A) \cap C_{n-1}(B) = C_{n-1}(A \cap B)$, it is a cycle, and
\[
\partial[z] = [\partial a] \in H_{n-1}(A \cap B).
\]
Indeed, in the description of the connecting map in Theorem~\ref{homological-algebra:theorem-long-exact-sequence}, $(a, b)$ is a preimage of $z$, and its boundary
$(\partial a, \partial b) = (\partial a, -\partial a)$ is the image of $\partial a$. In words: cut the cycle into a piece in $A$ and a piece in $B$; the common boundary of the
pieces is the image of the cycle.
\end{remark}

\section{Spheres, good pairs and degree}\label{singular-homology:section-spheres}

\begin{theorem}\label{singular-homology:theorem-spheres}
For $n \geq 0$, $\tilde{H}_k(S^n) \cong \ZZ$ for $k = n$ and $0$ otherwise. For $n \geq 1$ the Mayer--Vietoris connecting map is an isomorphism
$\tilde{H}_k(S^n) \to \tilde{H}_{k-1}(S^{n-1})$, where $S^{n-1} \subseteq S^n$ is the equator $\{x_{n+1} = 0\}$.
\end{theorem}

\begin{proof}
For $n = 0$, $S^0$ is two points and $\tilde{H}_0(S^0) \cong \ZZ$, generated by the difference of the two points. For $n \geq 1$ let
$A = \{x \in S^n : x_{n+1} > -1/2\}$ and $B = \{x_{n+1} < 1/2\}$. By stereographic projection from $\mp e_{n+1}$ (Example~\ref{general-topology:example-stereographic}), $A$ and
$B$ are homeomorphic to open balls, hence contractible, and $A \cap B$ deformation retracts onto the equator via
$H(x, t) = ((1 - t)x + t(x', 0))/\|(1 - t)x + t(x', 0)\|$, where $x = (x', x_{n+1})$ with $x' \neq 0$. The reduced Mayer--Vietoris sequence gives the
isomorphism, and induction gives the result.
\end{proof}

\begin{example}\label{singular-homology:example-circle-generator}
\emph{A generator of $H_1(S^1)$.} For $n = 1$ the proof uses $A = \{x_2 > -\frac12\}$ and $B = \{x_2 < \frac12\}$ in $S^1 \subseteq \RR^2$; their intersection consists of two
arcs, one containing $p = (1, 0)$ and one containing $q = (-1, 0)$. Let
\[
a(t) = \bigl(1 - 2t, 2\sqrt{t(1 - t)}\bigr), \qquad b(t) = \bigl(2t - 1, -2\sqrt{t(1 - t)}\bigr), \qquad t \in I,
\]
the upper half circle from $p$ to $q$ and the lower half circle from $q$ to $p$ (indeed $(1 - 2t)^2 + 4t(1 - t) = 1$). Then $a$ lies in $A$, $b$ lies in $B$, and $z = a + b$
is a cycle, as in Example~\ref{singular-homology:example-cycles}(1). By Remark~\ref{singular-homology:remark-mayer-vietoris-boundary}, $\partial[z] = [\partial a] = [\{q\} - \{p\}]$,
which generates $\tilde{H}_0(A \cap B) \cong \ZZ$, as $p$ and $q$ lie in different path components. Since $\partial$ is an isomorphism, $[z]$ generates $H_1(S^1) \cong \ZZ$. In
particular $z$ is not a boundary: the circle has a hole that the cycle $z$ goes around once. Neither half $a$ nor $b$ is a cycle by itself; the connecting map
detects precisely how the two halves are glued.
\end{example}

\begin{definition}\label{singular-homology:definition-good-pair}
A pair $(X, A)$ is \emph{good} if $A$ is a nonempty closed subspace that is a deformation retract of some open neighbourhood $V$ of $A$ in $X$.
\end{definition}

\begin{proposition}\label{singular-homology:proposition-good-pair}
For a good pair, the quotient map $q \colon (X, A) \to (X/A, A/A)$ induces isomorphisms $H_n(X, A) \cong H_n(X/A, A/A) \cong \tilde{H}_n(X/A)$. For example, for $n \geq 1$,
$(D^n, S^{n-1})$ is good and $H_k(D^n, S^{n-1}) \cong \tilde{H}_k(S^n) \cong \tilde{H}_{k-1}(S^{n-1})$.
\end{proposition}

\begin{proof}
Consider the commutative diagram
\begin{center}
\begin{tikzcd}
H_n(X, A) \arrow[r] \arrow[d, "q_*"] & H_n(X, V) \arrow[d, "q_*"] & H_n(X \setminus A, V \setminus A) \arrow[l] \arrow[d, "q_*"] \\
H_n(X/A, A/A) \arrow[r] & H_n(X/A, V/A) & H_n(X/A \setminus A/A, V/A \setminus A/A) \arrow[l]
\end{tikzcd}
\end{center}
The upper left map is an isomorphism by the long exact sequence of the triple $(X, V, A)$, since $H_n(V, A) = 0$ (Example~\ref{singular-homology:example-pair-point}). The
upper right map is an isomorphism by excision, as $\overline{A} = A \subseteq V = \operatorname{int}V$. Since $V$ is open and saturated, the deformation retraction of $V$
induces one of $V/A$ onto the point $A/A$ (Lemma~\ref{homotopy:lemma-quotient-times-interval}), so the lower maps are isomorphisms for the same reasons. The right
vertical map is induced by a homeomorphism of pairs, because $q$ restricts to a bijective quotient map $X \setminus A \to X/A \setminus A/A$ (again by
Lemma~\ref{homotopy:lemma-quotient-times-interval}, as $X \setminus A$ is open and saturated). So the left vertical map is an isomorphism. The last isomorphism is
Example~\ref{singular-homology:example-pair-point}. For $(D^n, S^{n-1})$ take $V = D^n \setminus \{0\}$ (Example~\ref{homotopy:example-contractible}); $D^n/S^{n-1}$ is
compact, and the map $D^n \to S^n$, $x \mapsto (2\sqrt{1 - \|x\|^2}\,x, 2\|x\|^2 - 1)$, induces a continuous bijection $D^n/S^{n-1} \to S^n$, hence a homeomorphism. The
last isomorphism follows from the long exact sequence of $(D^n, S^{n-1})$.
\end{proof}

\begin{definition}\label{singular-homology:definition-degree}
The \emph{degree} of a map $f \colon S^n \to S^n$ is the integer $\deg f$ such that $f_*$ is multiplication by $\deg f$ on $\tilde{H}_n(S^n) \cong \ZZ$.
\end{definition}

\begin{proposition}\label{singular-homology:proposition-degree}
\begin{enumerate}
\item $\deg \id = 1$, $\deg(fg) = \deg f \deg g$, and homotopic maps have the same degree; a map that is not surjective has degree $0$.
\item A reflection $x \mapsto (x_1, \ldots, -x_i, \ldots, x_{n+1})$ has degree $-1$, and the antipodal map has degree $(-1)^{n+1}$.
\item A map $f \colon S^n \to S^n$ without fixed points is homotopic to the antipodal map, so $\deg f = (-1)^{n+1}$.
\end{enumerate}
\end{proposition}

\begin{proof}
(1) Functoriality and homotopy invariance; a non-surjective map factors through $S^n \setminus \{\mathrm{pt}\} \cong \RR^n$, which is contractible.

(2) Let $r$ be the reflection in the first coordinate. For $n = 0$ it swaps the two points and acts by $-1$ on $\tilde{H}_0(S^0)$. For $n \geq 1$, $r$ preserves the sets $A$, $B$ of
the proof of Theorem~\ref{singular-homology:theorem-spheres} and commutes with the deformation retraction of $A \cap B$ onto the equator, on which it is again the
reflection in the first coordinate. By naturality of the Mayer--Vietoris sequence, $\deg r$ on $S^n$ equals $\deg r$ on $S^{n-1}$, hence equals $-1$. A reflection in the
$i$-th coordinate is $\pi r \pi^{-1}$ for a coordinate permutation $\pi$, so has degree $-1$ by (1). The antipodal map is the composite of $n + 1$ reflections.

(3) $(1 - t)f(x) - tx$ never vanishes: otherwise $f(x) = \frac{t}{1-t}x$ with $\|f(x)\| = \|x\| = 1$ forces $f(x) = x$. Normalising gives a homotopy from $f$ to $-\id$.
\end{proof}

\begin{corollary}\label{singular-homology:corollary-sphere-applications}
\begin{enumerate}
\item (Brouwer) There is no retraction $D^{n+1} \to S^n$, and every map $D^{n+1} \to D^{n+1}$ has a fixed point.
\item (Hairy ball theorem) $S^n$ has a continuous nowhere vanishing tangent vector field ($v \colon S^n \to \RR^{n+1}$ with $\langle v(x), x \rangle = 0$) if and only if
$n$ is odd.
\item (Invariance of dimension) If nonempty open subsets $U \subseteq \RR^m$ and $V \subseteq \RR^n$ are homeomorphic, then $m = n$.
\end{enumerate}
\end{corollary}

\begin{proof}
(1) $\tilde{H}_n(S^n) \neq 0$ while $D^{n+1}$ is contractible, so the identity of $\tilde{H}_n(S^n)$ cannot factor through $\tilde{H}_n(D^{n+1}) = 0$. The formula in the proof
of Corollary~\ref{homotopy:corollary-brouwer} turns a fixed point free map into a retraction.

(2) If $n = 2k - 1$, $v(x) = (-x_2, x_1, \ldots, -x_{2k}, x_{2k-1})$ works. Conversely, if $v$ never vanishes, let $u = v/\|v\|$. Since $x \perp u(x)$ are unit vectors,
$(1 - t)x + tu(x)$ and $(1 - t)u(x) - tx$ never vanish, and normalising gives homotopies $\id \simeq u \simeq -\id$. So $1 = (-1)^{n+1}$ and $n$ is odd.

(3) Let $x \in U$. By excision (removing the closed set $\RR^m \setminus U$ from $\RR^m \setminus \{x\}$) and the long exact sequence,
$H_k(U, U \setminus \{x\}) \cong H_k(\RR^m, \RR^m \setminus \{x\}) \cong \tilde{H}_{k-1}(\RR^m \setminus \{x\}) \cong \tilde{H}_{k-1}(S^{m-1})$ for $m \geq 1$, which is nonzero exactly
for $k = m$. A homeomorphism $U \to V$ gives isomorphisms of these local homology groups. For $m = 0$ or $n = 0$ the spaces are points, and a nonempty open subset of
$\RR^n$ is a point only if $n = 0$.
\end{proof}

\section{Cellular homology}\label{singular-homology:section-cellular}

\begin{lemma}\label{singular-homology:lemma-cw-homology}
Let $X$ be a CW complex (Definition~\ref{homotopy:definition-cw}).
\begin{enumerate}
\item $H_k(X^n, X^{n-1})$ is zero for $k \neq n$, and free abelian with basis the classes $e_\alpha$ of the characteristic maps $\Phi_\alpha \colon (D^n, S^{n-1}) \to (X^n, X^{n-1})$
applied to a fixed generator of $H_n(D^n, S^{n-1}) \cong \ZZ$, for $k = n$.
\item $H_k(X^n) = 0$ for $k > n$.
\item The inclusion $X^n \to X$ induces an isomorphism $H_k(X^n) \to H_k(X)$ for $k < n$.
\end{enumerate}
\end{lemma}

\begin{proof}
(1) $X^{n-1}$ is closed in $X^n$ (as in the proof of Proposition~\ref{homotopy:proposition-attaching-cells}), and it is a deformation retract of the open set obtained by
removing the centres of the $n$-cells, via the radial deformation of $D^n \setminus \{0\}$ onto $S^{n-1}$, which descends to the quotient by
Lemma~\ref{homotopy:lemma-quotient-times-interval}. So $(X^n, X^{n-1})$ is a good pair, and so is $(\coprod_\alpha D^n_\alpha, \coprod_\alpha S^{n-1}_\alpha)$. The
characteristic maps induce a continuous bijection $\coprod_\alpha D^n_\alpha/\coprod_\alpha S^{n-1}_\alpha \to X^n/X^{n-1}$, which is a homeomorphism: a set whose
preimage in $\coprod_\alpha D^n_\alpha$ is open has open preimage in $X^{n-1} \sqcup \coprod_\alpha D^n_\alpha$, since that preimage meets $X^{n-1}$ in $\emptyset$ or
$X^{n-1}$. By Proposition~\ref{singular-homology:proposition-good-pair} and its naturality,
$H_k(X^n, X^{n-1}) \cong H_k(\coprod_\alpha D^n_\alpha, \coprod_\alpha S^{n-1}_\alpha) = \bigoplus_\alpha H_k(D^n, S^{n-1})$, which is $\bigoplus_\alpha \ZZ$ for $k = n$ and $0$
otherwise. (For $n = 0$, $X^{-1} = \emptyset$ and $H_k(X^0)$ is computed by Proposition~\ref{singular-homology:proposition-h0-point}.)

(2) By the exact sequence $H_k(X^{n-1}) \to H_k(X^n) \to H_k(X^n, X^{n-1})$ and induction on $n$, starting from $X^{-1} = \emptyset$.

(3) For $k < n$ the exact sequence $H_{k+1}(X^{n+1}, X^n) \to H_k(X^n) \to H_k(X^{n+1}) \to H_k(X^{n+1}, X^n)$ has outer terms zero, so $H_k(X^n) \cong H_k(X^m)$ for all
$m \geq n$. A singular chain has compact image, which lies in a finite subcomplex (Theorem~\ref{homotopy:theorem-cw-properties}(2)), hence in some $X^m$. So
every cycle of $X$ is a cycle of some $X^m$, and a cycle of $X^n$ that bounds in $X$ bounds in some $X^m$; this gives surjectivity and injectivity.
\end{proof}

\begin{definition}\label{singular-homology:definition-cellular}
The \emph{cellular chain complex} of a CW complex $X$ has $C^{CW}_n(X) = H_n(X^n, X^{n-1})$ and differential $d_n$ the composite
$H_n(X^n, X^{n-1}) \xrightarrow{\partial} H_{n-1}(X^{n-1}) \xrightarrow{j_{n-1}} H_{n-1}(X^{n-1}, X^{n-2})$. Then $d_{n-1}d_n = 0$ because it contains two consecutive maps
$\partial \circ j_{n-1}$ of the long exact sequence of $(X^{n-1}, X^{n-2})$. Its homology is the \emph{cellular homology} $H^{CW}_n(X)$.
\end{definition}

\begin{theorem}\label{singular-homology:theorem-cellular}
$H^{CW}_n(X) \cong H_n(X)$ for every CW complex $X$.
\end{theorem}

\begin{proof}
Write $j_n \colon H_n(X^n) \to H_n(X^n, X^{n-1})$, which is injective since $H_n(X^{n-1}) = 0$ by Lemma~\ref{singular-homology:lemma-cw-homology}(2), and
$\partial_{n+1} \colon H_{n+1}(X^{n+1}, X^n) \to H_n(X^n)$, so $d_{n+1} = j_n\partial_{n+1}$. Since $j_n$ is injective, $\ker d_n = \ker(j_{n-1}\partial_n) = \ker\partial_n = \im j_n$, and
$j_n$ maps $\im\partial_{n+1}$ isomorphically onto $\im d_{n+1}$. Hence
$H^{CW}_n(X) \cong \im j_n/\im d_{n+1} \cong H_n(X^n)/\im\partial_{n+1}$. By the exact sequence $H_{n+1}(X^{n+1}, X^n) \xrightarrow{\partial_{n+1}} H_n(X^n) \to H_n(X^{n+1}) \to H_n(X^{n+1}, X^n) = 0$,
this is $H_n(X^{n+1})$, which is $H_n(X)$ by Lemma~\ref{singular-homology:lemma-cw-homology}(3).
\end{proof}

\begin{theorem}[Cellular boundary formula]\label{singular-homology:theorem-cellular-boundary}
For $n \geq 2$, $d_n(e^n_\alpha) = \sum_\beta d_{\alpha\beta}e^{n-1}_\beta$, where $d_{\alpha\beta}$ is the degree of the composite
$S^{n-1} \xrightarrow{\varphi_\alpha} X^{n-1} \to X^{n-1}/(X^{n-1} \setminus e^{n-1}_\beta) \cong S^{n-1}$ of the attaching map with the map collapsing the complement of
the cell $e^{n-1}_\beta$ (identified with $S^{n-1}$ via $\Phi_\beta$). For $n = 1$, $d_1(e^1_\alpha)$ is the difference of the endpoint and the starting point of the
$1$-cell.
\end{theorem}

This is \cite[Section 2.2, Cellular Boundary Formula]{hatcher-at}; the proof is a diagram chase using naturality of the long exact sequences for the characteristic
maps.

\begin{example}\label{singular-homology:example-cellular-computations}
\begin{enumerate}
\item If all cells of $X$ are in even dimensions, all $d_n$ vanish and $H_n(X)$ is free with basis the $n$-cells. For $\CC\PP^n$ (one cell in each even dimension
$\leq 2n$), $H_k(\CC\PP^n) \cong \ZZ$ for $k$ even with $0 \leq k \leq 2n$, and $0$ otherwise.
\item For the closed orientable surface $\Sigma_g$ of genus $g$ (Example~\ref{homotopy:example-cw}), $d_1 = 0$ (one $0$-cell) and $d_2 = 0$ (each $1$-cell appears twice
in the attaching word with opposite signs), so $H_0 = \ZZ$, $H_1 = \ZZ^{2g}$, $H_2 = \ZZ$.
\item For $\RR\PP^n$ with one cell in each dimension, $d_k$ is multiplication by $1 + (-1)^k$ \cite[Example 2.42]{hatcher-at}. So $H_k(\RR\PP^n)$ is $\ZZ$ for $k = 0$,
$\ZZ/2\ZZ$ for $k$ odd with $0 < k < n$, $\ZZ$ for $k = n$ odd, and $0$ otherwise.
\end{enumerate}
\end{example}

\begin{definition}\label{singular-homology:definition-euler-characteristic}
The \emph{Euler characteristic} of a finite CW complex with $c_n$ cells of dimension $n$ is $\chi(X) = \sum_n (-1)^nc_n$.
\end{definition}

\begin{proposition}\label{singular-homology:proposition-euler-characteristic}
$\chi(X) = \sum_n (-1)^n\dim_\QQ(H_n(X) \otimes \QQ)$. In particular $\chi(X)$ depends only on the homotopy type of $X$.
\end{proposition}

\begin{proof}
$\QQ = S^{-1}\ZZ$ is a flat $\ZZ$-module (Proposition~\ref{commutative-algebra:proposition-localisation}), so $H_n(C^{CW} \otimes \QQ) \cong H_n(C^{CW}) \otimes \QQ \cong H_n(X) \otimes \QQ$.
For a bounded complex $V$ of finite-dimensional vector spaces, rank--nullity (Theorem~\ref{linear-algebra:theorem-rank-nullity}) gives
$\dim V_n = \dim Z_n + \dim B_{n-1}$ and $\dim H_n = \dim Z_n - \dim B_n$, so $\sum(-1)^n\dim V_n = \sum(-1)^n\dim H_n$. Apply this to $C^{CW} \otimes \QQ$, with $\dim = c_n$.
\end{proof}

\begin{example}\label{singular-homology:example-euler}
$\chi(S^n) = 1 + (-1)^n$, $\chi(\Sigma_g) = 2 - 2g$, $\chi(\CC\PP^n) = n + 1$, $\chi(\RR\PP^n) = 1$ for $n$ even and $0$ for $n$ odd. As $\chi(\Sigma_g)$ determines $g$, surfaces of
different genus are not homotopy equivalent.
\end{example}

\section{Homology with coefficients}\label{singular-homology:section-coefficients}

\begin{definition}\label{singular-homology:definition-coefficients}
For an abelian group $G$, $C_n(X, A; G) = C_n(X, A) \otimes_\ZZ G$, and $H_n(X, A; G)$ is its homology. For a ring $R$ these are $R$-modules.
\end{definition}

\begin{theorem}\label{singular-homology:theorem-coefficients}
\begin{enumerate}
\item Homotopy invariance, the long exact sequences, excision and Mayer--Vietoris hold for $H_\bullet(-; G)$, and cellular homology with coefficients, computed from
$C^{CW}_\bullet(X) \otimes G$, agrees with $H_\bullet(X; G)$.
\item (Universal coefficients) There is a natural short exact sequence
$0 \to H_n(X, A) \otimes G \to H_n(X, A; G) \to \Tor_1^\ZZ(H_{n-1}(X, A), G) \to 0$. It splits, non-naturally \cite[Theorem 3A.3]{hatcher-at}.
\end{enumerate}
\end{theorem}

\begin{proof}
(1) The complexes $C(X, A)$ are free, so the short exact sequences of chain complexes used above are degreewise split and remain exact after $\otimes G$; chain homotopies
remain chain homotopies. The arguments therefore apply verbatim, and the proof of Theorem~\ref{singular-homology:theorem-cellular} goes through with
$H_n(X^n, X^{n-1}; G) \cong C^{CW}_n(X) \otimes G$. (2) Apply the K\"unneth formula (Theorem~\ref{homological-algebra:theorem-kunneth}) to the free complex $C(X, A)$ and
the complex $G$ concentrated in degree $0$. The splitting is quoted.
\end{proof}

\begin{example}\label{singular-homology:example-rp-mod-2}
With coefficients in $\mathbb{F}_2 = \ZZ/2\ZZ$ all cellular differentials of $\RR\PP^n$ vanish, since $1 + (-1)^k$ is even; so $H_k(\RR\PP^n; \mathbb{F}_2) \cong \mathbb{F}_2$ for
$0 \leq k \leq n$. With coefficients in a field $F$ of characteristic $0$, $H_\bullet(X; F) \cong H_\bullet(X) \otimes F$, because $F$ is a $\QQ$-vector space, hence a direct sum of copies of the flat module $\QQ$, so $\Tor_1^\ZZ(-, F) = 0$
(Proposition~\ref{homological-algebra:proposition-flat-tor}).
\end{example}

\section{The Hurewicz theorems}\label{singular-homology:section-hurewicz}

\begin{theorem}[Hurewicz in degree one]\label{singular-homology:theorem-hurewicz}
Let $X$ be path-connected and $x_0 \in X$. Regarding a loop as a singular $1$-simplex via $\Delta^1 \cong I$, $(t_0, t_1) \mapsto t_1$, gives a well-defined homomorphism
$h \colon \pi_1(X, x_0) \to H_1(X)$, which induces an isomorphism $\pi_1(X, x_0)^{\mathrm{ab}} \cong H_1(X)$ from the abelianisation
(Proposition~\ref{groups:proposition-abelianization}).
\end{theorem}

\begin{proof}
For paths we write $f \sim g$ if $f - g \in C_1(X)$ is a boundary. (i) A constant path $c$ satisfies $c \sim 0$: it is the boundary of the constant $2$-simplex.
(ii) If $f \simeq g$ relative to endpoints via $F \colon I^2 \to X$, divide $I^2$ into two triangles along the diagonal and compose with $F$; the boundary of the
difference of the two resulting $2$-simplices is $f - g$ plus constant paths, so $f \sim g$. (iii) $f * g \sim f + g$: the $2$-simplex $\sigma = (f * g) \circ \pi$, where
$\pi \colon \Delta^2 \to I$ is the affine map with $e_0 \mapsto 0$, $e_1 \mapsto 1/2$, $e_2 \mapsto 1$, has $\partial\sigma = g - f * g + f$. (iv) Hence
$f + \bar f \sim f * \bar f \sim c \sim 0$. So $h$ is well defined and a homomorphism, and it factors through $\pi_1^{\mathrm{ab}}$ as $H_1(X)$ is abelian.

Choose for each $x \in X$ a path $\gamma_x$ from $x_0$ to $x$, with $\gamma_{x_0}$ constant. For a singular $1$-simplex $\tau$ let $L(\tau) = [\gamma_{\tau(0)} * \tau * \bar\gamma_{\tau(1)}]$
in $\pi_1(X, x_0)^{\mathrm{ab}}$, written additively, and extend $L$ linearly to $C_1(X)$. For a singular $2$-simplex $\sigma$ with edges $\tau_0, \tau_1, \tau_2$ (so
$\partial\sigma = \tau_0 - \tau_1 + \tau_2$), $L(\partial\sigma)$ is the class of the loop $\gamma_{\sigma(e_0)} * \tau_2 * \tau_0 * \bar\tau_1 * \bar\gamma_{\sigma(e_0)}$, which is
null-homotopic because the loop around the boundary of $\Delta^2$ is null-homotopic in the contractible space $\Delta^2$. So $L$ induces
$\bar L \colon H_1(X) \to \pi_1(X, x_0)^{\mathrm{ab}}$. Clearly $\bar L h = \id$. For a cycle $z = \sum n_i\tau_i$, by (iii) and (iv),
$hL(z) = \sum n_i(\gamma_{\tau_i(0)} + \tau_i - \gamma_{\tau_i(1)}) = z - \Gamma(\partial z) = z$ in $H_1(X)$, where $\Gamma \colon C_0(X) \to C_1(X)$ is the homomorphism $x \mapsto \gamma_x$.
So $h$ and $\bar L$ are inverse isomorphisms.
\end{proof}

In higher degrees, homotopy groups are related to homology in the same way, as long as the lower homotopy groups vanish. Let $n \geq 1$ and fix a generator $[S^n]$ of
$H_n(S^n) \cong \ZZ$ (Theorem~\ref{singular-homology:theorem-spheres}). A map $f \colon (I^n, \partial I^n) \to (X, x_0)$ factors through the quotient $I^n/\partial I^n \cong S^n$, and so
induces $f_* \colon H_n(S^n) \to H_n(X)$. The \emph{Hurewicz map} is
\[
h_n \colon \pi_n(X, x_0) \to H_n(X), \qquad [f] \mapsto f_*[S^n] .
\]
It is well defined, since homotopic maps induce the same map in homology (Theorem~\ref{singular-homology:theorem-homotopy-invariance}), and it is a homomorphism \cite[Section 4.2]{hatcher-at}; for $n = 1$ it is the map $h$ of
Theorem~\ref{singular-homology:theorem-hurewicz}. A space $X$ is \emph{$n$-connected} if it is path-connected and $\pi_k(X, x_0) = 0$ for $1 \leq k \leq n$; so $0$-connected means
path-connected and $1$-connected means simply connected.

\begin{theorem}[Hurewicz]\label{singular-homology:theorem-hurewicz-general}
Let $n \geq 2$ and let $X$ be an $(n-1)$-connected space. Then $\tilde{H}_k(X) = 0$ for $k < n$, and the Hurewicz map $h_n \colon \pi_n(X, x_0) \to H_n(X)$ is an isomorphism.
\end{theorem}

This is \cite[Theorem 4.32]{hatcher-at}, where a relative version is also proved; we do not prove it here. In words: \emph{the first nonvanishing homotopy group of a simply connected space
is its first nonvanishing reduced homology group.} For $n = 1$ the analogous statement is Theorem~\ref{singular-homology:theorem-hurewicz}, with the abelianisation; for $n \geq 2$ the group
$\pi_n$ is already abelian (Proposition~\ref{homotopy:proposition-higher-abelian}), and no correction is needed.

\begin{example}\label{singular-homology:example-hurewicz}
\begin{enumerate}
\item For $n \geq 2$ the sphere $S^n$ is $(n-1)$-connected (Remark~\ref{homotopy:remark-higher-computations}), so $\pi_n(S^n) \cong H_n(S^n) \cong \ZZ$, generated by the class of the
identity. Composing with $h_n$, the homotopy class of a map $S^n \to S^n$ is detected by its degree.
\item A wedge $\bigvee_{\alpha}S^n$ of $n$-spheres with $n \geq 2$ is $(n-1)$-connected (by cellular approximation, as for a single sphere), so
$\pi_n(\bigvee_\alpha S^n) \cong \bigoplus_\alpha\ZZ$ by Exercise~\ref{singular-homology:exercise-wedge}: one generator for each sphere.
\item The Hurewicz theorem says nothing about the groups above the first nonvanishing one. For $S^2$, $\pi_3(S^2) \cong \ZZ$ while $H_3(S^2) = 0$
(Remark~\ref{homotopy:remark-higher-computations}): higher homotopy groups are not determined by homology.
\end{enumerate}
\end{example}

\begin{counterexample}\label{singular-homology:counterexample-hurewicz-simply-connected}
For $n = 2$ the hypothesis that $X$ is simply connected cannot be dropped. Let $X = S^1 \vee S^2$. Then $H_2(X) \cong \ZZ$ (Exercise~\ref{singular-homology:exercise-wedge}), but $\pi_2(X)$ is a
free abelian group of infinite rank. Indeed, the universal cover $\widetilde{X}$ is the real line with a copy of $S^2$ attached at each integer, and by
Proposition~\ref{homotopy:proposition-covering-higher} $\pi_2(X) \cong \pi_2(\widetilde{X})$. Collapsing the contractible line shows that $\widetilde{X}$ is homotopy equivalent to the
wedge $\bigvee_{m \in \ZZ}S^2$, which is simply connected, so $\pi_2(\widetilde{X}) \cong \bigoplus_{m \in \ZZ}\ZZ$ by Example~\ref{singular-homology:example-hurewicz}(2). The group
$\pi_1(X) \cong \ZZ$ acts on $\pi_2(X)$ by translating the spheres, and $H_2(X)$ only sees the quotient by this action.
\end{counterexample}

\begin{remark}\label{singular-homology:remark-homology-whitehead}
Combined with Whitehead's theorem (Theorem~\ref{homotopy:theorem-whitehead}), the Hurewicz theorem gives a criterion that only involves homology: a map between simply connected CW complexes that
induces isomorphisms on all homology groups is a homotopy equivalence \cite[Corollary 4.33]{hatcher-at}. Simple connectivity is needed here too: the Poincar\'e homology sphere with a point
removed has the homology of a point, but it is not contractible, since its fundamental group is a nontrivial perfect group (equal to its commutator subgroup, so that its abelianisation $H_1$
vanishes).
\end{remark}

\begin{remark}\label{singular-homology:remark-eilenberg-steenrod}
Homotopy invariance, the long exact sequence, excision, additivity over disjoint unions and the value on a point are the \emph{Eilenberg--Steenrod axioms}. On
CW complexes they determine homology, as the proof of Theorem~\ref{singular-homology:theorem-cellular} shows. Other constructions satisfying them, such as
simplicial, \v{C}ech or sheaf-theoretic homology on suitable spaces, agree with singular homology there; the comparison with de Rham cohomology is
Chapter~\ref{de-rham}.
\end{remark}

\section{Singular Cohomology}\label{singular-homology:section-singular-cohomology}
Cohomology is obtained by dualising singular chains. As groups it carries the same information as homology, by the universal
coefficient theorem, but it has a natural ring structure given by the cup product, which distinguishes spaces with the same
homology (Example~\ref{singular-homology:example-cp2-wedge}). Cohomology classes are the objects of the Hodge conjecture: a
Hodge class is a rational cohomology class of a special type, and the conjecture asks whether it is a combination of classes of
subvarieties. The cap product relates cohomology and homology and leads to Poincar\'e duality (Chapter~\ref{poincare-duality}).
We follow \cite[Chapter 3]{hatcher-at}.

\noindent\textbf{Prerequisites.} Chapter~\ref{singular-homology} (Singular Homology).

Throughout, $R$ is a commutative ring and $G$ an abelian group, written additively.

\section{Cochains and cohomology}\label{singular-homology:section-cochains}

\begin{definition}\label{singular-homology:definition-cohomology}
For a pair $(X, A)$ let $C^n(X, A; G) = \Hom_\ZZ(C_n(X, A), G)$, the group of \emph{singular $n$-cochains}, with coboundary
$\delta \colon C^n \to C^{n+1}$, $\delta\varphi = \varphi \circ \partial$. So $C^n(X, A; G)$ is the group of functions from singular $n$-simplices of $X$ to $G$ vanishing on
simplices in $A$, and $(\delta\varphi)(\sigma) = \sum_i(-1)^i\varphi(\sigma \circ \delta^i)$. The \emph{singular cohomology} $H^n(X, A; G)$ is the cohomology of this complex,
and $H^n(X; G) = H^n(X, \emptyset; G)$. A map of pairs $f$ induces $f^\# = \Hom(f_\#, G)$ and $f^* \colon H^n(Y, B; G) \to H^n(X, A; G)$, contravariantly. Reduced cohomology
$\tilde{H}^n(X; G)$ is obtained by dualising the augmented chain complex. We write $H^n(X)$ for $H^n(X; \ZZ)$. The \emph{Kronecker pairing}
$\langle -, - \rangle \colon H^n(X, A; G) \times H_n(X, A) \to G$ is $\langle[\varphi], [c]\rangle = \varphi(c)$.
\end{definition}

\begin{theorem}\label{singular-homology:theorem-axioms}
\begin{enumerate}
\item Homotopic maps of pairs induce the same map on cohomology.
\item For $B \subseteq A \subseteq X$ there is a natural long exact sequence
$\cdots \to H^n(X, A; G) \to H^n(X, B; G) \to H^n(A, B; G) \xrightarrow{\delta} H^{n+1}(X, A; G) \to \cdots$, and
$\langle \delta\alpha, c \rangle = \langle \alpha, \partial c \rangle$ for $\alpha \in H^n(A; G)$ and $c \in H_{n+1}(X, A)$.
\item (Excision) If $Z \subseteq A \subseteq X$ with $\overline{Z} \subseteq \operatorname{int}A$, then $H^n(X, A; G) \to H^n(X \setminus Z, A \setminus Z; G)$ is an isomorphism.
\item (Mayer--Vietoris) If $\operatorname{int}A \cup \operatorname{int}B = X$, there is a long exact sequence
$\cdots \to H^n(X; G) \to H^n(A; G) \oplus H^n(B; G) \to H^n(A \cap B; G) \to H^{n+1}(X; G) \to \cdots$.
\item $H^0(X; G)$ is the group of functions from the set of path components of $X$ to $G$.
\end{enumerate}
\end{theorem}

\begin{proof}
Applying $\Hom(-, G)$ turns chain homotopies into cochain homotopies, which gives (1) from Theorem~\ref{singular-homology:theorem-homotopy-invariance}. The
sequences $0 \to C_n(A, B) \to C_n(X, B) \to C_n(X, A) \to 0$ and $0 \to C_n(A \cap B) \to C_n(A) \oplus C_n(B) \to C^{\{A, B\}}_n(X) \to 0$ consist of free abelian groups
and split degreewise, since the quotients are free on the simplices not in the subspace. So they remain exact after $\Hom(-, G)$, and
Theorem~\ref{homological-algebra:theorem-long-exact-sequence} gives the long exact sequences. The connecting map sends $[\varphi]$ to the class of $\delta\tilde\varphi$ for
an extension $\tilde\varphi$ of $\varphi$ to $X$, which gives the formula for the pairing. The chain homotopy equivalence $C^{\mathcal{U}}(X) \to C(X)$ of
Proposition~\ref{singular-homology:proposition-small-chains} dualises to a cochain homotopy equivalence, which gives (3) as in
Theorem~\ref{singular-homology:theorem-excision} and allows the replacement of $H^n(\Hom(C^{\{A,B\}}(X), G))$ by $H^n(X; G)$ in (4). For (5), a $0$-cocycle is a function
$\varphi$ on points with $\varphi(\gamma(1)) - \varphi(\gamma(0)) = 0$ for every path $\gamma$.
\end{proof}

\begin{theorem}[Universal coefficients]\label{singular-homology:theorem-universal-coefficients}
There is a natural short exact sequence
\[
0 \to \Ext^1_\ZZ(H_{n-1}(X, A), G) \to H^n(X, A; G) \to \Hom(H_n(X, A), G) \to 0,
\]
where the second map is $\alpha \mapsto \langle \alpha, - \rangle$. It splits, non-naturally. If $G = F$ is a field and $A = \emptyset$, then
$H^n(X; F) \cong \Hom_F(H_n(X; F), F)$.
\end{theorem}

\begin{proof}
The first statement is Theorem~\ref{homological-algebra:theorem-universal-coefficients} for the free complex $C(X, A)$. For a field,
$\Hom_\ZZ(C_n(X), F) = \Hom_F(C_n(X) \otimes F, F)$, and $\Hom_F(-, F)$ is exact, so it commutes with taking cohomology.
\end{proof}

\begin{lemma}\label{singular-homology:lemma-ext-z}
$\Ext^1_\ZZ(\ZZ, G) = 0$ and $\Ext^1_\ZZ(\ZZ/m\ZZ, G) \cong G/mG$; $\Ext^1_\ZZ$ commutes with finite direct sums. Hence, if $H_{n-1}(X)$ and $H_n(X)$ are finitely generated,
$H^n(X; \ZZ) \cong \ZZ^{b_n} \oplus T_{n-1}$, where $b_n$ is the rank of $H_n(X)$ and $T_{n-1}$ the torsion subgroup of $H_{n-1}(X)$.
\end{lemma}

\begin{proof}
Use the projective resolution $0 \to \ZZ \xrightarrow{m} \ZZ \to \ZZ/m\ZZ \to 0$ and Theorem~\ref{homological-algebra:theorem-balance}: $\Ext^1(\ZZ/m\ZZ, G)$ is the cokernel of
$G \xrightarrow{m} G$. Additivity follows from additivity of $\Hom$ and of derived functors. For the last statement, write $H_n(X) \cong \ZZ^{b_n} \oplus T_n$ and $H_{n-1}(X) \cong \ZZ^{b_{n-1}} \oplus T_{n-1}$ with $T_n$, $T_{n-1}$ finite direct sums of groups
$\ZZ/m\ZZ$ (Theorem~\ref{groups:theorem-structure-abelian}). By the split universal coefficient sequence,
$H^n(X) \cong \Hom(H_n(X), \ZZ) \oplus \Ext^1(H_{n-1}(X), \ZZ)$; here $\Hom(\ZZ/m\ZZ, \ZZ) = 0$ gives $\Hom(H_n(X), \ZZ) \cong \ZZ^{b_n}$, and
$\Ext^1(\ZZ/m\ZZ, \ZZ) \cong \ZZ/m\ZZ$ gives $\Ext^1(H_{n-1}(X), \ZZ) \cong T_{n-1}$.
\end{proof}

\begin{example}\label{singular-homology:example-cohomology-groups}
\begin{enumerate}
\item $\tilde{H}^k(S^n; G) \cong G$ for $k = n$ and $0$ otherwise.
\item By Example~\ref{singular-homology:example-cellular-computations} and Lemma~\ref{singular-homology:lemma-ext-z}, $H^k(\RR\PP^n; \ZZ)$ is $\ZZ$ for $k = 0$, $\ZZ/2\ZZ$ for
$k$ even with $0 < k \leq n$, $\ZZ$ for $k = n$ odd, and $0$ otherwise: torsion moves up one degree compared with homology. With coefficients in $\mathbb{F}_2$,
$H^k(\RR\PP^n; \mathbb{F}_2) \cong \mathbb{F}_2$ for $0 \leq k \leq n$.
\end{enumerate}
\end{example}

\begin{theorem}[Cellular cohomology]\label{singular-homology:theorem-cellular-singular-cohomology}
For a CW complex $X$, $H^n(X; G)$ is isomorphic to the cohomology of the complex $\Hom(C^{CW}_\bullet(X), G)$ with the transposed differentials.
\end{theorem}

\begin{proof}
By the universal coefficient theorem and Lemma~\ref{singular-homology:lemma-cw-homology}: $H^k(X^n, X^{n-1}; G)$ is $\Hom(C^{CW}_n(X), G)$ for $k = n$ and $0$ otherwise;
$H^k(X^n; G) = 0$ for $k > n$, because $H_k(X^n) = 0$ for $k > n$ and $H_n(X^n)$ is a subgroup of the free group $C^{CW}_n(X)$, hence free
(Proposition~\ref{homological-algebra:proposition-pid-submodules-free}). From the long exact sequence of $(X, X^n)$ and
Lemma~\ref{singular-homology:lemma-cw-homology}(3), $H_k(X, X^n) = 0$ for $k \leq n$, hence $H^k(X, X^n; G) = 0$ for $k \leq n$ and $H^k(X; G) \to H^k(X^n; G)$ is an
isomorphism for $k < n$. The proof of Theorem~\ref{singular-homology:theorem-cellular} now applies with all arrows reversed. The coboundary
$H^n(X^n, X^{n-1}) \to H^n(X^n) \to H^{n+1}(X^{n+1}, X^n)$ is the transpose of $d_{n+1}$ by the formula $\langle\delta\alpha, c\rangle = \langle\alpha, \partial c\rangle$.
\end{proof}

\begin{corollary}\label{singular-homology:corollary-even-cells}
Let $X$ be a CW complex with finitely many cells in each dimension and no cells of odd dimension. Then $H^k(X; G) = 0$ for $k$ odd, and $H^k(X; G) \cong G^{c_k}$ for $k$ even, where $c_k$ is the
number of $k$-cells. The same holds for homology, and $\chi(X) = \sum_kc_k$ when $X$ is finite.
\end{corollary}

\begin{proof}
The cellular chain complex has $C^{CW}_k(X) = 0$ for $k$ odd, so all its differentials vanish, being maps into or out of zero groups. By Theorem~\ref{singular-homology:theorem-cellular-singular-cohomology} the cohomology is
$\Hom(C^{CW}_k(X), G) = \Hom(\ZZ^{c_k}, G) = G^{c_k}$ in each even degree $k$, and $0$ in odd degrees; the homology statement is Theorem~\ref{singular-homology:theorem-cellular}. The Euler characteristic is
the alternating sum of the ranks.
\end{proof}

\begin{example}\label{singular-homology:example-even-cells}
$\CC\PP^n$ has one cell in each even dimension $0, 2, \ldots, 2n$ (Example~\ref{singular-homology:example-cellular-computations}), so $H^{2k}(\CC\PP^n; \ZZ) \cong \ZZ$ for $0 \leq k \leq n$ and all other groups
vanish, and $\chi(\CC\PP^n) = n + 1$. The product $S^2 \times S^2$, with cells $e^0$, $e^2 \times e^0$, $e^0 \times e^2$, $e^2 \times e^2$, has $H^0 = H^4 = \ZZ$ and $H^2 = \ZZ^2$. For $\RR\PP^n$, which has one
cell in each dimension $\leq n$, the corollary does not apply, and indeed its cohomology has torsion (Example~\ref{singular-homology:example-cohomology-groups}).
\end{example}

\section{The cup product}\label{singular-homology:section-cup}

\begin{definition}\label{singular-homology:definition-cup}
For $\varphi \in C^k(X; R)$ and $\psi \in C^l(X; R)$, the \emph{cup product} $\varphi \smile \psi \in C^{k+l}(X; R)$ is
\[
(\varphi \smile \psi)(\sigma) = \varphi(\sigma \circ [e_0, \ldots, e_k]) \cdot \psi(\sigma \circ [e_k, \ldots, e_{k+l}]),
\]
the product being taken in $R$: the ``front $k$-face'' times the ``back $l$-face''. We write $\sigma|[e_0, \ldots, e_k]$ for $\sigma \circ [e_0, \ldots, e_k]$.
\end{definition}

\begin{lemma}\label{singular-homology:lemma-cup-coboundary}
$\delta(\varphi \smile \psi) = \delta\varphi \smile \psi + (-1)^k\varphi \smile \delta\psi$. The cup product of cochains is associative, $R$-bilinear, has unit the $0$-cochain $1$
with value $1$ on every point, and is natural: $f^\#(\varphi \smile \psi) = f^\#\varphi \smile f^\#\psi$.
\end{lemma}

\begin{proof}
Let $\sigma \colon \Delta^{k+l+1} \to X$ and $n = k + l + 1$. Then
\[
(\delta\varphi \smile \psi)(\sigma) = \sum_{i=0}^{k+1}(-1)^i\varphi(\sigma|[e_0, \ldots, \widehat{e_i}, \ldots, e_{k+1}])\psi(\sigma|[e_{k+1}, \ldots, e_n]),
\]
\[
(-1)^k(\varphi \smile \delta\psi)(\sigma) = \sum_{i=k}^{n}(-1)^i\varphi(\sigma|[e_0, \ldots, e_k])\psi(\sigma|[e_k, \ldots, \widehat{e_i}, \ldots, e_n]).
\]
The term $i = k + 1$ of the first sum cancels the term $i = k$ of the second. The remaining terms are $\sum_i (-1)^i(\varphi \smile \psi)(\sigma|[e_0, \ldots, \widehat{e_i}, \ldots, e_n])$,
which is $(\varphi \smile \psi)(\partial\sigma) = \delta(\varphi \smile \psi)(\sigma)$. The other statements are immediate from the definition.
\end{proof}

\begin{definition}\label{singular-homology:definition-cohomology-ring}
By Lemma~\ref{singular-homology:lemma-cup-coboundary}, the cup product of two cocycles is a cocycle, and the product of a cocycle and a coboundary is a coboundary. It
therefore induces $\smile \colon H^k(X; R) \times H^l(X; R) \to H^{k+l}(X; R)$, making $H^*(X; R) = \bigoplus_k H^k(X; R)$ an associative graded $R$-algebra with unit, the
\emph{cohomology ring}, and $f^*$ a homomorphism of graded rings. If $\varphi$ vanishes on chains in $A$ so does $\varphi \smile \psi$, which gives a product
$H^k(X, A; R) \times H^l(X; R) \to H^{k+l}(X, A; R)$.
\end{definition}

\begin{theorem}[Graded commutativity]\label{singular-homology:theorem-graded-commutative}
For $\alpha \in H^k(X; R)$ and $\beta \in H^l(X; R)$, $\alpha \smile \beta = (-1)^{kl}\beta \smile \alpha$.
\end{theorem}

\begin{proof}[Sketch of proof]
Let $\bar\sigma = \sigma \circ [e_n, \ldots, e_0]$ be $\sigma$ with its vertices in reverse order, and $\rho(\sigma) = \varepsilon_n\bar\sigma$ with $\varepsilon_n = (-1)^{n(n+1)/2}$.
A direct computation shows that $\rho$ is a chain map and that
$\rho^\#\varphi \smile \rho^\#\psi = (-1)^{kl}\rho^\#(\psi \smile \varphi)$, using $\varepsilon_{k+l} = (-1)^{kl}\varepsilon_k\varepsilon_l$. The omitted verification is that $\rho$ is
chain homotopic to the identity, by the homotopy $P(\sigma) = \sum_i (-1)^i\varepsilon_{n-i}\,\sigma \circ \pi \circ [v_0, \ldots, v_i, w_n, \ldots, w_i]$, where
$v_j = (e_j, 0)$ and $w_j = (e_j, 1)$ are the vertices of $\Delta^n \times I$ and $\pi \colon \Delta^n \times I \to \Delta^n$ is the projection; see
\cite[Theorem 3.14]{hatcher-at}. Then $\rho^* = \id$ on cohomology, and the formula follows.
\end{proof}

\begin{example}\label{singular-homology:example-wedge-ring}
For a wedge $X \vee Y$ of CW complexes (at $0$-cells), the inclusions induce an isomorphism of rings
$\tilde{H}^*(X \vee Y; R) \cong \tilde{H}^*(X; R) \times \tilde{H}^*(Y; R)$, where products between the two factors vanish: the isomorphism of groups follows from
Exercise~\ref{singular-homology:exercise-wedge} and the universal coefficient theorem, and it is induced by the restrictions to $X$ and $Y$, which are ring homomorphisms.
In particular all cup products $\tilde{H}^k(S^m) \times \tilde{H}^l(S^m) \to \tilde{H}^{k+l}(S^m)$ vanish, for degree reasons.
\end{example}

\section{The K\"unneth formula}\label{singular-homology:section-kunneth}

\begin{definition}\label{singular-homology:definition-cross-product}
The \emph{cross product} of $a \in H^k(X; R)$ and $b \in H^l(Y; R)$ is $a \times b = p_X^*(a) \smile p_Y^*(b) \in H^{k+l}(X \times Y; R)$, where $p_X$, $p_Y$ are the
projections. It is $R$-bilinear and induces $H^*(X; R) \otimes_R H^*(Y; R) \to H^*(X \times Y; R)$, a ring homomorphism for the product
$(a \otimes b)(c \otimes d) = (-1)^{|b||c|}ac \otimes bd$ on the tensor product.
\end{definition}

\begin{theorem}[K\"unneth formula]\label{singular-homology:theorem-kunneth}
Let $X$ and $Y$ be CW complexes such that $H^k(Y; R)$ is a finitely generated free $R$-module for every $k$ (for instance $R$ a field and $Y$ a finite complex).
Then the cross product $H^*(X; R) \otimes_R H^*(Y; R) \to H^*(X \times Y; R)$ is an isomorphism of graded rings.
\end{theorem}

This is \cite[Theorem 3.15]{hatcher-at}. The ring homomorphism property follows from graded commutativity:
$(a \times b)(c \times d) = p_X^*a\, p_Y^*b\, p_X^*c\, p_Y^*d = (-1)^{|b||c|}p_X^*(ac)\,p_Y^*(bd)$. The homology K\"unneth formula for fields, and the
Eilenberg--Zilber theorem $C_\bullet(X \times Y) \simeq C_\bullet(X) \otimes C_\bullet(Y)$ underlying it, are in \cite[Section 3.B]{hatcher-at}.

\begin{example}\label{singular-homology:example-torus}
$H^*(S^1; \ZZ) = \Lambda_\ZZ[x]$ with $|x| = 1$ (an exterior algebra, as $x^2 \in H^2(S^1) = 0$). By induction and the K\"unneth formula, the cohomology ring of the torus
$T^n = (S^1)^n$ is the exterior algebra $\Lambda_\ZZ[x_1, \ldots, x_n]$ with $|x_i| = 1$, so $H^k(T^n) \cong \ZZ^{\binom{n}{k}}$. Similarly
$H^*(S^2 \times S^2; \ZZ) = \ZZ[a, b]/(a^2, b^2)$ with $|a| = |b| = 2$.
\end{example}

\section{Cohomology rings of projective spaces}\label{singular-homology:section-projective}

\begin{theorem}\label{singular-homology:theorem-projective-spaces}
$H^*(\RR\PP^n; \mathbb{F}_2) \cong \mathbb{F}_2[\alpha]/(\alpha^{n+1})$ with $|\alpha| = 1$, and $H^*(\CC\PP^n; \ZZ) \cong \ZZ[\alpha]/(\alpha^{n+1})$ with $|\alpha| = 2$. Also
$H^*(\RR\PP^\infty; \mathbb{F}_2) = \mathbb{F}_2[\alpha]$ and $H^*(\CC\PP^\infty; \ZZ) = \ZZ[\alpha]$.
\end{theorem}

This is \cite[Theorem 3.12]{hatcher-at}; the additive structure was computed above, and the content is that the generator in each degree is a power of $\alpha$. A
proof via Poincar\'e duality is given in Chapter~\ref{poincare-duality}.

\begin{example}\label{singular-homology:example-cp2-wedge}
$\CC\PP^2$ and $S^2 \vee S^4$ have the same homology and cohomology groups ($\ZZ$ in degrees $0, 2, 4$), but they are not homotopy equivalent: in $H^*(\CC\PP^2)$ the square of
a generator of $H^2$ generates $H^4$, whereas in $H^*(S^2 \vee S^4)$ it is zero by Example~\ref{singular-homology:example-wedge-ring}. Likewise $S^2 \times S^2$ and
$S^2 \vee S^2 \vee S^4$ have the same groups but different rings.
\end{example}

\begin{corollary}[Borsuk--Ulam]\label{singular-homology:corollary-borsuk-ulam}
There is no map $g \colon S^n \to S^{n-1}$ with $g(-x) = -g(x)$ for all $x$. Consequently every map $f \colon S^n \to \RR^n$ identifies some pair of antipodal points.
\end{corollary}

\begin{proof}
For $n = 1$ an odd map $S^1 \to S^0$ cannot exist, since its image would be a single point. Let $n \geq 2$ and suppose $g$ exists; it induces
$\bar{g} \colon \RR\PP^n \to \RR\PP^{n-1}$. For a path-connected space $X$, the universal coefficient theorem and Theorem~\ref{singular-homology:theorem-hurewicz} give
$H^1(X; \mathbb{F}_2) \cong \Hom(H_1(X), \mathbb{F}_2) \cong \Hom(\pi_1(X), \mathbb{F}_2)$, naturally. For $m \geq 1$ let $w_m \colon \pi_1(\RR\PP^m) \to \mathbb{F}_2$ be the monodromy of
the two-sheeted covering $S^m \to \RR\PP^m$ (Example~\ref{homotopy:example-coverings}): $w_m[\gamma] = 0$ if and only if the lifts of $\gamma$ are loops. It is a
homomorphism (Corollary~\ref{homotopy:corollary-lifting-consequences}(3)) and nonzero. Since $\pi_1(\RR\PP^n) \cong \ZZ/2\ZZ$, $w_n$ corresponds to the generator $\alpha$ of
$H^1(\RR\PP^n; \mathbb{F}_2)$, and $\pi_1(\RR\PP^n)$ is generated by the image of a path $\gamma$ in $S^n$ from $x$ to $-x$. Its image under $\bar{g}$ is the image of the path
$g \circ \gamma$ from $g(x)$ to $-g(x)$, whose lift is not a loop; so $w_{n-1}(\bar{g}_*[\gamma]) = 1 = w_n[\gamma]$ and $\bar{g}^*(w_{n-1}) = \alpha$. Then
$\bar{g}^*(w_{n-1}^n) = \alpha^n \neq 0$ by Theorem~\ref{singular-homology:theorem-projective-spaces}, but $w_{n-1}^n \in H^n(\RR\PP^{n-1}; \mathbb{F}_2) = 0$, a contradiction.
If $f(x) \neq f(-x)$ for all $x$, then $g(x) = (f(x) - f(-x))/\|f(x) - f(-x)\|$ is an odd map $S^n \to S^{n-1}$.
\end{proof}

\section{Cup length and coverings by contractible sets}\label{singular-homology:section-cup-length}

\begin{lemma}\label{singular-homology:lemma-relative-cup}
Let $A$, $B$ be open subsets of $X$. The cup product of cochains induces a product $H^k(X, A; R) \times H^l(X, B; R) \to H^{k+l}(X, A \cup B; R)$, compatible with the absolute cup product under the maps
$H^\bullet(X, C; R) \to H^\bullet(X; R)$ for $C = A, B, A \cup B$.
\end{lemma}

\begin{proof}
If $\varphi$ vanishes on simplices in $A$ and $\psi$ on simplices in $B$, then $\varphi \smile \psi$ vanishes on simplices in $A$ (because of its front face) and on simplices in $B$ (because of its back face),
hence on the subcomplex $C(A) + C(B) \subseteq C(A \cup B)$. By Lemma~\ref{singular-homology:lemma-cup-coboundary}, this gives a product with values in the cohomology of
$\Hom(C(X)/(C(A) + C(B)), R)$. It remains to identify this cohomology with $H^\bullet(X, A \cup B; R)$. The inclusion $C(A) + C(B) \to C(A \cup B)$ is a chain homotopy equivalence by
Proposition~\ref{singular-homology:proposition-small-chains}, applied to the open cover $\{A, B\}$ of $A \cup B$, so it induces isomorphisms on cohomology with coefficients in $R$. The sequences
$0 \to C(A) + C(B) \to C(X) \to C(X)/(C(A) + C(B)) \to 0$ and $0 \to C(A \cup B) \to C(X) \to C(X, A \cup B) \to 0$ split degreewise, the quotients being free on the simplices not in the subcomplex, so
they give long exact cohomology sequences after applying $\Hom(-, R)$, and the inclusions give a map between them. The five lemma (Lemma~\ref{homological-algebra:lemma-five}) shows that
$H^\bullet(X, A \cup B; R) \to H^\bullet(\Hom(C(X)/(C(A) + C(B)), R))$ is an isomorphism. Compatibility with the absolute product holds on cochains.
\end{proof}

\begin{proposition}\label{singular-homology:proposition-lusternik-schnirelmann}
Let $X$ be covered by open sets $U_1, \ldots, U_n$ such that each inclusion $U_i \to X$ is homotopic to a constant map. Then $\alpha_1 \smile \cdots \smile \alpha_n = 0$ for all
$\alpha_i \in H^{k_i}(X; R)$ with $k_i \geq 1$.
\end{proposition}

\begin{proof}
The restriction $H^{k}(X; R) \to H^{k}(U_i; R)$ is induced by the inclusion, which is homotopic to a map factoring through a point, so it is zero for $k \geq 1$, as $H^k(\mathrm{pt}; R) = 0$
(Theorem~\ref{singular-homology:theorem-axioms}(1)). By the long exact sequence of the pair $(X, U_i)$, each $\alpha_i$ lifts to $\tilde\alpha_i \in H^{k_i}(X, U_i; R)$. By
Lemma~\ref{singular-homology:lemma-relative-cup}, applied $n - 1$ times, $\tilde\alpha_1 \smile \cdots \smile \tilde\alpha_n$ lies in $H^\bullet(X, U_1 \cup \cdots \cup U_n; R) = H^\bullet(X, X; R) = 0$ and
maps to $\alpha_1 \smile \cdots \smile \alpha_n$.
\end{proof}

\begin{corollary}\label{singular-homology:corollary-cup-length}
\begin{enumerate}
\item The torus $T^n$ cannot be covered by $n$ open sets that are contractible in $T^n$, but it can be covered by $n + 1$ of them.
\item $\CC\PP^n$ cannot be covered by $n$ open sets contractible in $\CC\PP^n$, and it is covered by the $n + 1$ standard affine charts, each homeomorphic to $\CC^n$.
\item $\RR\PP^n$ cannot be covered by $n$ open sets contractible in $\RR\PP^n$.
\end{enumerate}
\end{corollary}

\begin{proof}
By Proposition~\ref{singular-homology:proposition-lusternik-schnirelmann} it suffices to exhibit $n$ classes of positive degree with nonzero product: $x_1 \smile \cdots \smile x_n$ generates $H^n(T^n)$
(Example~\ref{singular-homology:example-torus}), and $\alpha^n \neq 0$ in $H^{2n}(\CC\PP^n; \ZZ)$ and in $H^n(\RR\PP^n; \mathbb{F}_2)$ (Theorem~\ref{singular-homology:theorem-projective-spaces}). For the
covering of $T^n = \RR^n/\ZZ^n$ by $n + 1$ sets, let $U_j$, $0 \leq j \leq n$, be the image of the open set of points $x \in \RR^n$ whose coordinates avoid the interval
$[\frac{j}{n+1} + \frac{1}{4(n+1)}, \frac{j}{n+1} + \frac{3}{4(n+1)}]$ modulo $1$: each coordinate then ranges in an open interval of $\RR/\ZZ$, so $U_j$ is homeomorphic to an open cube, which is contractible.
Every $x$ lies in some $U_j$, because its $n$ coordinates can meet at most $n$ of the $n + 1$ disjoint intervals.
\end{proof}

\begin{remark}\label{singular-homology:remark-ls-category}
The least $m$ such that $X$ is covered by $m + 1$ open sets contractible in $X$ is the \emph{Lusternik--Schnirelmann category} of $X$. Proposition~\ref{singular-homology:proposition-lusternik-schnirelmann} bounds it from
below by the \emph{cup length}, the largest number of classes of positive degree with nonzero product. By Corollary~\ref{singular-homology:corollary-cup-length} the category of $T^n$ is $n$. Lusternik and
Schnirelmann introduced this invariant to bound from below the number of critical points of smooth functions on closed manifolds.
\end{remark}

\section{The cap product}\label{singular-homology:section-cap}

\begin{definition}\label{singular-homology:definition-cap}
For $\sigma \colon \Delta^n \to X$ and $\varphi \in C^k(X; R)$ with $k \leq n$, the \emph{cap product} is
\[
\sigma \frown \varphi = \varphi(\sigma|[e_0, \ldots, e_k]) \cdot \sigma|[e_k, \ldots, e_n] \in C_{n-k}(X; R),
\]
extended $R$-linearly to $C_n(X; R) = C_n(X) \otimes R$.
\end{definition}

\begin{proposition}\label{singular-homology:proposition-cap}
\begin{enumerate}
\item $\partial(\sigma \frown \varphi) = (-1)^k(\partial\sigma \frown \varphi - \sigma \frown \delta\varphi)$, so the cap product induces
$\frown \colon H_n(X; R) \times H^k(X; R) \to H_{n-k}(X; R)$, and relative versions $H_n(X, A; R) \times H^k(X; R) \to H_{n-k}(X, A; R)$ and
$H_n(X, A; R) \times H^k(X, A; R) \to H_{n-k}(X; R)$.
\item (Naturality) $f_*(c) \frown \varphi = f_*(c \frown f^*\varphi)$ for $f \colon X \to Y$, $c \in H_n(X; R)$ and $\varphi \in H^k(Y; R)$.
\item $\psi(c \frown \varphi) = (\varphi \smile \psi)(c)$ for $\psi \in C^{n-k}(X; R)$, and $(c \frown \varphi) \frown \psi = c \frown (\varphi \smile \psi)$.
\end{enumerate}
\end{proposition}

\begin{proof}
(1) We have
$\partial(\sigma \frown \varphi) = \sum_{i \geq k}(-1)^{i-k}\varphi(\sigma|[e_0, \ldots, e_k])\,\sigma|[e_k, \ldots, \widehat{e_i}, \ldots, e_n]$, while
\[
\partial\sigma \frown \varphi = \sum_{i \leq k}(-1)^i\varphi(\sigma|[e_0, \ldots, \widehat{e_i}, \ldots, e_{k+1}])\,\sigma|[e_{k+1}, \ldots, e_n]
+ \sum_{i \geq k+1}(-1)^i\varphi(\sigma|[e_0, \ldots, e_k])\,\sigma|[e_k, \ldots, \widehat{e_i}, \ldots, e_n]
\]
and $\sigma \frown \delta\varphi = \sum_{i \leq k+1}(-1)^i\varphi(\sigma|[e_0, \ldots, \widehat{e_i}, \ldots, e_{k+1}])\,\sigma|[e_{k+1}, \ldots, e_n]$. Subtracting, what remains is the
second sum together with the term $-(-1)^{k+1}\varphi(\sigma|[e_0, \ldots, e_k])\,\sigma|[e_{k+1}, \ldots, e_n]$, which is the term $i = k$; this is
$(-1)^k\partial(\sigma \frown \varphi)$. Hence cycle $\frown$ cocycle is a cycle, boundary $\frown$ cocycle and cycle $\frown$ coboundary are boundaries. For the relative
versions, $\sigma \frown \varphi$ lies in $C(A)$ if $\sigma$ does, and vanishes if $\varphi$ vanishes on $A$ and $\sigma$ lies in $A$. (2) and (3) are immediate from the
definitions: $\psi(\sigma \frown \varphi) = \varphi(\sigma|[e_0, \ldots, e_k])\psi(\sigma|[e_k, \ldots, e_n]) = (\varphi \smile \psi)(\sigma)$.
\end{proof}

\begin{remark}\label{singular-homology:remark-duality-preview}
For a closed oriented $n$-manifold $M$ with fundamental class $[M] \in H_n(M; \ZZ)$, the map $H^k(M; R) \to H_{n-k}(M; R)$, $\alpha \mapsto [M] \frown \alpha$, is an
isomorphism: this is Poincar\'e duality (Chapter~\ref{poincare-duality}). Combined with Proposition~\ref{singular-homology:proposition-cap}(3) it shows that the
cup product pairing $H^k(M; F) \times H^{n-k}(M; F) \to F$, $(\alpha, \beta) \mapsto \langle \alpha \smile \beta, [M]\rangle$, is nondegenerate over a field $F$.
\end{remark}

\section{Exercises}\label{singular-homology:section-exercises}

\subsection{Singular Homology}\label{singular-homology:subsection-exercises-singular-homology}

\begin{exercise}\label{singular-homology:exercise-spheres-not-equivalent}
Show that $S^m$ and $S^n$ are homotopy equivalent only if $m = n$, and that $\RR^m \setminus \{0\}$ and $\RR^n \setminus \{0\}$ are homeomorphic only if $m = n$.
\end{exercise}

\begin{solution}
By Theorem~\ref{singular-homology:theorem-spheres}, $\tilde{H}_\bullet(S^m)$ is nonzero only in degree $m$. The spaces $\RR^m \setminus \{0\}$ and $S^{m-1}$ are homotopy equivalent
(Example~\ref{homotopy:example-contractible}).
\end{solution}

\begin{exercise}\label{singular-homology:exercise-suspension}
The \emph{suspension} $SX$ is the quotient of $X \times [-1, 1]$ collapsing $X \times \{1\}$ and $X \times \{-1\}$ to two points. Show that $\tilde{H}_{n+1}(SX) \cong \tilde{H}_n(X)$ for
all $n$.
\end{exercise}

\begin{solution}
Let $A$ and $B$ be the images of $X \times (-1/2, 1]$ and $X \times [-1, 1/2)$. They are contractible (slide to the cone points; the homotopies descend by
Lemma~\ref{homotopy:lemma-quotient-times-interval}), and $A \cap B \cong X \times (-1/2, 1/2) \simeq X$. The reduced Mayer--Vietoris sequence gives the isomorphism.
\end{solution}

\begin{exercise}\label{singular-homology:exercise-wedge}
Show that $\tilde{H}_n(\bigvee_\alpha X_\alpha) \cong \bigoplus_\alpha\tilde{H}_n(X_\alpha)$ if the basepoints $x_\alpha$ are such that each $(X_\alpha, x_\alpha)$ is a good pair.
Compute $H_\bullet$ of a wedge of $k$ circles.
\end{exercise}

\begin{solution}
The pair $(\coprod X_\alpha, \coprod\{x_\alpha\})$ is good with quotient $\bigvee X_\alpha$, so
$\tilde{H}_n(\bigvee X_\alpha) \cong H_n(\coprod X_\alpha, \coprod\{x_\alpha\}) = \bigoplus H_n(X_\alpha, x_\alpha) = \bigoplus\tilde{H}_n(X_\alpha)$. For $k$ circles, $H_0 = \ZZ$, $H_1 = \ZZ^k$,
and $H_n = 0$ for $n \geq 2$; this agrees with Theorem~\ref{singular-homology:theorem-hurewicz}, since the abelianisation of a free group of rank $k$ is $\ZZ^k$.
\end{solution}

\begin{exercise}\label{singular-homology:exercise-degree-comparison}
Using Theorem~\ref{singular-homology:theorem-hurewicz} and Remark~\ref{homotopy:remark-exponential}, show that the degree of a map $S^1 \to S^1$ defined via $\pi_1$
agrees with the homological degree.
\end{exercise}

\begin{exercise}\label{singular-homology:exercise-lefschetz}
Let $X$ be a finite CW complex and $f \colon X \to X$ a cellular map. Define the Lefschetz number $\Lambda(f) = \sum_n(-1)^n\operatorname{tr}(f_* \mid H_n(X; \QQ))$. Show that
$\Lambda(f) = \sum_n (-1)^n \operatorname{tr}(f_\# \mid C^{CW}_n(X) \otimes \QQ)$, and deduce that $\Lambda(\id) = \chi(X)$.
\end{exercise}

\begin{solution}
A cellular map sends $X^n$ into $X^n$, so it is a map of pairs $(X^n, X^{n-1}) \to (X^n, X^{n-1})$ and induces maps $f_\#$ on $C^{CW}_n(X)$; they commute with the
differentials, which are composites of maps of long exact sequences. The isomorphism $H^{CW}_n(X) \cong H_n(X)$ of Theorem~\ref{singular-homology:theorem-cellular} is built
from maps induced by inclusions and connecting maps, so it is compatible with $f$; the same holds with coefficients in $\QQ$ (Theorem~\ref{singular-homology:theorem-coefficients}).

It remains to prove the \emph{Hopf trace formula}: if $V$ is a bounded complex of finite-dimensional vector spaces and $\varphi$ a chain map from $V$ to itself, then
$\sum_n(-1)^n\operatorname{tr}(\varphi \mid V_n) = \sum_n(-1)^n\operatorname{tr}(\varphi \mid H_n(V))$. The subspaces $B_n \subseteq Z_n \subseteq V_n$ are preserved by $\varphi$.
Choosing a basis of $B_n$, extending it to $Z_n$ and then to $V_n$, the matrix of $\varphi$ on $V_n$ is block upper triangular, with diagonal blocks the matrices of $\varphi$
on $B_n$, on $Z_n/B_n = H_n$ and on $V_n/Z_n$. The differential induces an isomorphism $V_n/Z_n \cong B_{n-1}$ compatible with $\varphi$, so
$\operatorname{tr}(\varphi \mid V_n) = \operatorname{tr}(\varphi \mid B_n) + \operatorname{tr}(\varphi \mid H_n) + \operatorname{tr}(\varphi \mid B_{n-1})$. In the alternating sum the
terms with $B$ cancel, as in the proof of Proposition~\ref{singular-homology:proposition-euler-characteristic}. For $f = \id$, the trace on $C^{CW}_n(X) \otimes \QQ \cong \QQ^{c_n}$ is
$c_n$, so $\Lambda(\id) = \chi(X)$.
\end{solution}

\begin{exercise}\label{singular-homology:exercise-klein-bottle}
The Klein bottle $K$ is obtained from $S^1 \vee S^1$ by attaching a $2$-cell along $aba^{-1}b$. Compute $\pi_1(K)$, $H_\bullet(K)$ and $H_\bullet(K; \mathbb{F}_2)$.
\end{exercise}

\begin{solution}
By Proposition~\ref{homotopy:proposition-attaching-cells} and Example~\ref{homotopy:example-wedge}, $\pi_1(K) \cong \langle a, b \mid aba^{-1}b \rangle$. The cellular chain complex
has one $0$-cell, two $1$-cells $a, b$ and one $2$-cell $e$. Since there is one $0$-cell, $d_1 = 0$. The degrees in the cellular boundary formula are the sums of the
exponents of each letter in the attaching word, as in Example~\ref{singular-homology:example-cellular-computations}(2) \cite[Section 2.2]{hatcher-at}: $a$ occurs with
exponents $+1$ and $-1$, and $b$ twice with exponent $+1$. So $d_2(e) = 2b$, and
\[
H_0(K) \cong \ZZ, \qquad H_1(K) \cong \ZZ^2/\langle(0, 2)\rangle \cong \ZZ \oplus \ZZ/2\ZZ, \qquad H_2(K) = \ker d_2 = 0.
\]
This agrees with Theorem~\ref{singular-homology:theorem-hurewicz}: abelianising the relation $aba^{-1}b = 1$ gives $2b = 0$. With coefficients in $\mathbb{F}_2$ the differential
$d_2 = 2b$ becomes $0$, so $H_k(K; \mathbb{F}_2)$ is $\mathbb{F}_2$, $\mathbb{F}_2^2$, $\mathbb{F}_2$ for $k = 0, 1, 2$; for instance $H_2(K; \mathbb{F}_2) \cong \Tor_1^\ZZ(\ZZ/2\ZZ, \mathbb{F}_2) \cong \mathbb{F}_2$ by
universal coefficients. Both computations give $\chi(K) = 1 - 2 + 1 = 0$. Unlike the torus, $H_2(K; \ZZ) = 0$: the $2$-cell is not a cycle, which reflects the fact that $K$ is not
orientable.
\end{solution}

\begin{exercise}\label{singular-homology:exercise-free-actions}
Let $n$ be even.
\begin{enumerate}
\item Show that for every map $f \colon S^n \to S^n$ there is a point $x$ with $f(x) = x$ or $f(x) = -x$.
\item Show that if a group $G$ acts on $S^n$ by homeomorphisms \emph{freely}, that is, no $g \neq 1$ has a fixed point, then $G$ has at most two elements. (The
antipodal map gives a free action of $\ZZ/2\ZZ$, with quotient $\RR\PP^n$.)
\end{enumerate}
For odd $n$ both statements fail: on $S^1 \subseteq \CC$, the rotation by a right angle moves every point to a point that is neither equal nor opposite to it, and
rotations by multiples of $2\pi/k$ give a free action of $\ZZ/k\ZZ$.
\end{exercise}

\begin{solution}
(1) Let $\alpha$ be the antipodal map, of degree $(-1)^{n+1} = -1$ (Proposition~\ref{singular-homology:proposition-degree}(2)). If $f$ has no fixed point, then
$\deg f = -1$ by Proposition~\ref{singular-homology:proposition-degree}(3). If $f(x) \neq -x$ for all $x$, then $\alpha \circ f$ has no fixed point, so
$-\deg f = \deg(\alpha f) = -1$ and $\deg f = 1$. The two cannot both hold.

(2) Each $g \in G$ acts by a homeomorphism, and $\deg g \cdot \deg g^{-1} = \deg\id = 1$, so $\deg g = \pm1$, and $g \mapsto \deg g$ is a homomorphism $G \to \{\pm1\}$ by
Proposition~\ref{singular-homology:proposition-degree}(1). If $g \neq 1$, then $g$ has no fixed point, so $\deg g = -1$. Hence the homomorphism is injective, and $|G| \leq 2$.
\end{solution}

\subsection{Singular Cohomology}\label{singular-homology:subsection-exercises-singular-cohomology}

\begin{exercise}\label{singular-homology:exercise-uct-rp}
Compute $H^k(\RR\PP^n; \ZZ/4\ZZ)$ from the cellular cochain complex, and check the result with the universal coefficient theorem.
\end{exercise}

\begin{solution}
The cellular cochain complex is $\ZZ/4 \to \ZZ/4 \to \cdots$ with differential $\Hom(d_{k+1}, \ZZ/4)$, multiplication by $1 + (-1)^{k+1}$: zero out of even degrees and
multiplication by $2$ out of odd degrees. So $H^0 = \ZZ/4$; for $0 < k < n$, $H^k \cong \ZZ/2$ (for $k$ odd the kernel of $2$ is $2\ZZ/4$ and the image of the previous zero map is
$0$; for $k$ even the kernel is all of $\ZZ/4$ and the image is $2\ZZ/4$); and $H^n$ is $\ZZ/4$ for $n$ odd and $\ZZ/2$ for $n$ even. With universal coefficients, for
$0 < k < n$ odd: $\Hom(\ZZ/2, \ZZ/4) \oplus \Ext(0, \ZZ/4) = \ZZ/2$; for $k$ even: $\Hom(0, \ZZ/4) \oplus \Ext(\ZZ/2, \ZZ/4) = \ZZ/4/2\ZZ/4 = \ZZ/2$.
\end{solution}

\begin{exercise}\label{singular-homology:exercise-cup-torus}
Let $\Sigma_g$ be the closed orientable surface of genus $g$. Using the map $\Sigma_g \to \bigvee_{i=1}^g T^2$ collapsing $g - 1$ circles, show that $H^1(\Sigma_g) \cong \ZZ^{2g}$
has a basis $a_i, b_i$ with $a_i \smile b_i = -b_i \smile a_i$ a generator of $H^2(\Sigma_g) \cong \ZZ$, and all other products of basis elements zero.
\end{exercise}

\begin{exercise}\label{singular-homology:exercise-hopf-invariant}
For a map $f \colon S^{2n-1} \to S^n$ with $n \geq 2$ let $X_f = S^n \cup_f D^{2n}$. Show that $H^n(X_f) \cong \ZZ \cdot a$ and $H^{2n}(X_f) \cong \ZZ \cdot b$ and define the Hopf
invariant $h(f)$ by $a^2 = h(f)b$. Show that $h(f) = 0$ if $n$ is odd, and that $h$ of the Hopf map $S^3 \to S^2$ is $\pm 1$, using $X_f = \CC\PP^2$.
\end{exercise}

\begin{exercise}\label{singular-homology:exercise-retract-ring}
Show that $S^2 \vee S^4$ is not a retract of $\CC\PP^2$, even up to homotopy: there are no maps $i \colon S^2 \vee S^4 \to \CC\PP^2$ and $r \colon \CC\PP^2 \to S^2 \vee S^4$ with
$ri \simeq \id$. Show also that there is no map $\CC\PP^2 \to S^2$ inducing an isomorphism on $H^2$.
\end{exercise}

\begin{solution}
A map $f \colon \CC\PP^2 \to S^2$ with $f^*$ an isomorphism on $H^2$ would give $f^*(x)^2 = f^*(x^2) = 0$ for a generator $x$ of $H^2(S^2)$, while $f^*(x)$ generates $H^2(\CC\PP^2)$
and its square generates $H^4(\CC\PP^2)$. If $S^2 \vee S^4 \xrightarrow{i} \CC\PP^2 \xrightarrow{r} S^2 \vee S^4$ composed to the identity (up to homotopy), then $i^*$ would be
surjective, and $r^*$ injective; a generator $u \in H^2(S^2 \vee S^4)$ would satisfy $r^*(u)^2 = r^*(u^2) = 0$ with $r^*(u) = \pm$ a generator of $H^2(\CC\PP^2)$ (as $r^*$ is
injective on $H^2 \cong \ZZ$ and $i^*r^* = \id$ forces $r^*$ to be split injective), contradicting $\alpha^2 \neq 0$.
\end{solution}
